% ==================================================================
%  TESI MAGISTRALE - Mauro Segreti
%  Versione con il RESTYLE di Luca + tutte le CORREZIONI del PDF
%  ------------------------------------------------------------------
%  I box rossi "DA FARE" segnalano le note di Luca che richiedono
%  lavoro tuo (nuove analisi, nuovi capitoli, ecc.).
%  Per farli sparire prima della consegna: \showtodostrue -> false
% ==================================================================
\documentclass[11pt,a4paper,twoside,openright]{report}
\PassOptionsToPackage{export}{adjustbox}

% --- Codifica, font, lingua ---
\usepackage[utf8]{inputenc}
\usepackage{lmodern}   % font vettoriale: serve a microtype con T1
\usepackage[T1]{fontenc}
\usepackage[english]{babel}
\usepackage{microtype}

% --- Matematica ---
\usepackage{amsmath, amssymb, amsfonts, amsthm}
\usepackage{bm}

% --- Figure e tabelle ---
\usepackage{graphicx}
\usepackage{adjustbox}
\usepackage{float}
\usepackage{booktabs}
\usepackage{caption}
\usepackage{subcaption}
\usepackage{lettrine}
\usepackage{tabularx}
\usepackage{array}
\graphicspath{{figures/}{TESI/}{TESI/Figures/}}

% --- DIDASCALIE IN ROSSO ---------------------------------------------
% Per tornare al nero prima della consegna: commenta le due righe qui sotto.
%\captionsetup{font={small,color=red}, labelfont={small,bf,color=red}}
%\captionsetup[sub]{font={small,color=red}, labelfont={small,bf,color=red}}

% --- Layout di pagina ---
\usepackage{geometry}
\geometry{a4paper, top=3cm, bottom=3cm, left=3.5cm, right=2.5cm} % margini per rilegatura
\usepackage{emptypage}
\usepackage[pagestyles]{titlesec}   % NON caricare fancyhdr: sono incompatibili

% --- Colori ---
\usepackage{xcolor}
\definecolor{sapphire}{rgb}{0.03, 0.15, 0.4}

% --- Box colorati ---
\usepackage[most]{tcolorbox}

% --- Link, bibliografia, riferimenti incrociati ---
% Ordine obbligatorio: hyperref -> biblatex -> cleveref
\usepackage{hyperref}
\hypersetup{
    colorlinks=true,
    linkcolor=sapphire,
    citecolor=sapphire,
    urlcolor=sapphire
}
\usepackage[numbers,sort&compress]{natbib}
\usepackage{cleveref}

% --- Formattazione titoli dei capitoli ---
\titleformat{\chapter}[display]
  {\normalfont\bfseries\color{sapphire}}
  {\filleft\huge\chaptertitlename\ \thechapter}
  {20pt}
  {\titlerule\vspace{1ex}\raggedright\Huge}
  [\vspace{1ex}\titlerule]

% --- Intestazioni di pagina ---
\newpagestyle{main}[\normalsize\scshape]{
  \sethead[\thepage][][\chaptertitlename\ \thechapter.\ \chaptertitle]
          {\toptitlemarks\thesection\quad\sectiontitle}{}{\thepage}
  \setheadrule{0.4pt}
}
\pagestyle{main}

% ==================================================================
%  BOX PROMEMORIA (note di Luca da lavorare)
%  Metti \showtodosfalse per nasconderli tutti nella versione finale
% ==================================================================
\newif\ifshowtodos
\showtodofalse
\newcommand{\faExclamation}{}
\newtcolorbox{todobox}[1][Nota di Luca]{
  enhanced, breakable,
  colback=red!4!white,
  colframe=red!65!black,
  coltitle=white,
  fonttitle=\bfseries\sffamily\small,
  fontupper=\small\sffamily,
  title={\faExclamation DA FARE --- #1},
  left=3mm, right=3mm, top=2mm, bottom=2mm,
  boxrule=0.8pt, arc=2mm
}
\newenvironment{todo}[1][Nota di Luca]
  {\ifshowtodos\begin{todobox}[#1]\else\setbox0\vbox\bgroup\fi}
  {\ifshowtodos\end{todobox}\else\egroup\fi}

% ==================================================================
%  MACRO PERSONALI (sempre dopo tutti i \usepackage)
% ==================================================================
\newcommand{\ccterm}[1]{%
  \overline{q}\gamma^\mu\left\{ig_w\frac{#1}{\sqrt{2}}\,\frac{1-\gamma^5}{2}\right\}q}

\begin{document}
% ==================================================================
%  FRONTESPIZIO
% ==================================================================
\begin{titlepage}
    \centering
    \vspace*{1cm}

    % ---------- Tor Vergata ----------
    \includegraphics[width=\textwidth]{logo_tor_vergata.png}

    \vspace*{\stretch{1.3}}

    % ---------- Dipartimento ----------
    {\large\color{sapphire}\textls[180]{\textsc{Department of Physics}}\par}
    \vspace{0.45cm}
    {\large Master's Degree in Physics\par}
    \vspace{0.15cm}
    {\small\itshape Curriculum in Fundamental Interactions and Experimental Techniques\par}

    \vspace*{\stretch{1.7}}

    % ---------- Titolo ----------
    {\color{sapphire}\rule{\textwidth}{0.8pt}}\\[1.1cm]
    {\huge\bfseries\color{sapphire}
     Muon performance studies with\\[0.25cm]
     the ATLAS Muon Spectrometer in Run 3\par}
    \vspace{1.1cm}
    {\color{sapphire}\rule{\textwidth}{0.8pt}}

    \vspace*{\stretch{1.7}}


    % ---------- Relatori e candidato ----------
    \begin{minipage}[t]{0.48\textwidth}
        \begin{flushleft}\large
            \emph{Supervisors:}\\[0.15cm]
            Prof. Anna \textsc{Di Ciaccio}\\
            Dr. Marco \textsc{Vanadia}\\
            Dr. Luca \textsc{Martinelli}
        \end{flushleft}
    \end{minipage}
    \hfill
    \begin{minipage}[t]{0.44\textwidth}
        \begin{flushright}\large
            \emph{Candidate:}\\[0.15cm]
            Mauro \textsc{Segreti}
        \end{flushright}
    \end{minipage}

    \vspace*{\stretch{1}}

    % ---------- Data ----------
    {\large Academic Year 2025/2026\par}
    \vspace{0.2cm}
    {\small\color{sapphire}Rome, Italy\par}
\end{titlepage}
% ==================================================================
%  ABSTRACT
% ==================================================================
\begin{abstract}
This thesis presents muon momentum resolution and performance studies for the ATLAS detector, based on the \textit{pp} collision data sample produced at \(\sqrt{s} = 13.6\ \text{TeV}\) at the LHC during Run~3, corresponding to an integrated luminosity of \(332\ \text{fb}^{-1}\) \cite{LuminosityEra,ATLAS:LumiRun3}. The muon momentum resolution are measured using simulated samples of $Z\rightarrow\mu^+\mu^-$ and $Z'\rightarrow\mu^+\mu-$ events. "DA RIVEDERE"
\end{abstract}

%CORRETTO--------------------------------------------------
% CORREZIONE (p.3): 13 TeV -> 13.6 TeV, energia del Run 3.
% CORREZIONE (p.3): 139 fb^-1 e' la luminosita' del Run 2 (140 dopo il ricalcolo).
%   Nel Run 3 ATLAS ha registrato 332 fb^-1. Vedi:
%   https://atlas.web.cern.ch/Atlas/GROUPS/DATAPREPARATION/PublicPlots/2026/DataSummary/figs/intlumivstimeRun123.png
% CORREZIONE (p.3): "This paper" -> "This thesis" ("ah si? e' diventato un paper questa tesi? ahahah")

% ==================================================================
%  INDICE
% ==================================================================
\cleardoublepage
\tableofcontents


% ==================================================================
%  CAPITOLO 1 - INTRODUZIONE
% ==================================================================
\chapter{Introduction}

\lettrine[lines=3,loversize=0.1,lraise=0.05]{T}{} he third data-taking period of the Large Hadron Collider, known as Run 3, started on 5 June 2022, at the end of two-year period of Long Shutdown 2, during which both the accelerator and the experiments underwent an extensive upgrade programme. Protons have been colliding at a center of mass energy of \(\sqrt{s}=13.6\ TeV\), which is the highest energy ever reached in a collider and a little increase with respect to those of \(\sqrt{s}=13\ TeV\) of Run 2.  The increase in energy was accompanied by a rise in instantaneous luminosity, which has passed \(2\times10^{34}\ cm^{-2}s^{-1}\), bringing to almost 60 the number of proton-proton collisions per bunch crossing. This pile-up level led to a hard constraint on the triggers and detector sensitivity. Run 3 ended in 2026, with its last \(pp\) collisions in May and Pb-Pb collisions in June; then the LHC entered the Long Shutdown 3, in view of the High-Luminosity LHC. In 4 years of data-taking, the ATLAS experiment recorded \(332\ fb^{-1}\) \(pp\) collision data, double with respect to Run 2, and the biggest ever. This simulated data sample, together with the detector configuration described in Chapter 3, forms the basis of this work. The thesis is organised as follows. Chapter~2 outlines the Standard Model, with particular attention to the production and decay of the $Z$ boson into the muon final state on which the analysis is based. Chapter~3 introduces the ATLAS experiment, while Chapter~4 focuses on the Muon Spectrometer and the upgrades it has undergone. Chapter~5 describes the procedures adopted to reconstruct and identify the muons of interest, and Chapter~6 presents the analysis of the data and Monte Carlo samples. Chapter~7 summarises the results and draws the conclusions. "PER ADESSO"



\clearpage

% ==================================================================
%  CAPITOLO 2 - MODELLO STANDARD
% ==================================================================
\chapter{A Brief Overview of the Standard Model}

\lettrine[lines=3,loversize=0.1,lraise=0.05]{F}{} in dagli anni '30 si è cercato di costruire una teoria che collegasse la Meccanica Quantistica alla Meccanica Relativistica. Gli studi di numerosi fisici hanno portato ad una brillante intuizione: che tutto l'universo fosse composto da entità elementari, chiamate particelle, e che questo "Zoo" componesse tutta la materia (e antimateria) che conosciamo. La teoria più accreditata e ricca di evidenze sperimentali è il Modello Standard, formulato attraverso la Teoria Quantistica dei Campi tra gli anni '60 e '70. Tuttavia, sebbene racchiuda una descrizione quasi completa di tutto il cosmo, è appunto ancora incompleto: non include la gravità, né spiega la materia oscura e le masse dei neutrini (e molto altro) e non basterebbe una tesi per descriverlo tutto. In questo capitolo racconterò i momenti salienti di questa teoria, e il successo più noto che forse conosciamo tutti, l'evidenza e scoperta del bosone di Higgs.

\section{Introduction}

L'Universo come lo conosciamo è composto da due classi di particelle, distinte dal
valore dello spin: i fermioni, a spin semintero, che ne costituiscono la materia, e i
bosoni, a spin intero, fra cui si trovano i mediatori delle interazioni fondamentali.
Le interazioni note sono quattro: elettromagnetica, forte, debole e gravitazionale. I fermioni si
suddividono a loro volta in leptoni e quark, ciascuno con la propria antiparticella. I leptoni, riportati nella \cref{tab:sm_leptons}, si organizzano in tre famiglie, o
"generazioni", identiche per numeri quantici e distinte unicamente dalla massa. Come i quark, sono particelle elementari: non mostrano alcuna struttura interna, e i limiti sperimentali attuali ne collocano le dimensioni al di sotto di $10^{-18}$~m. La prima
generazione costituisce tutta la materia stabile, mentre le altre due decadono verso di
essa. I leptoni carichi sono soggetti alle interazioni elettromagnetica e debole,
mentre i neutrini interagiscono soltanto debolmente. L'elettrone è stabile; il muone
decade per interazione debole con una vita media di $2.2~\mathrm{\mu s}$, e analogo discorso per il tau, ma con una vita media di $2.9 \cdot 10^{-13}$~s.  I quark, riportati nella \cref{tab:quarks}, hanno tutti spin $1/2$ e carica frazionaria, $+2/3\,|e|$ oppure $-1/3\,|e|$. Esistono in sei sapori, distinti da numeri quantici interni: up ($u$), down ($d$), charm ($c$), strange ($s$), bottom ($b$) e top ($t$). I quark $u$ e $d$ sono i più leggeri e hanno masse confrontabili, tanto da essere raggruppati in un 
\begin{figure}[htpb]
    \centering
    \includegraphics[width=\textwidth]{TESI/particles.png}
    \caption{Schematic representation of the elementary particles of the Standard Model: quarks, leptons, gauge bosons and the Higgs boson. \cite{CERN-SM-chart}}
    \label{fig:particles}
\end{figure}
doppietto di isospin\footnote{L'isospin fu introdotto da Heisenberg per descrivere protone e neutrone come due stati dello stesso oggetto rispetto all'interazione forte; riflettendo la loro somiglianza direttamente nella simmetria di sapore fra $u$ e $d$.}  $I = 1/2$, con terza componente $I_3 = +1/2$ per $u$ e $I_3 = -1/2$ per $d$. Agli altri sapori è associato un numero quantico
dedicato: stranezza $S = -1$ per il quark $s$, charm $C = +1$ per il quark $c$,
beauty $B = -1$ per il quark $b$ e topness $T = +1$ per il quark $t$. Ad ogni quark corrisponde un antiquark con carica elettrica e numeri di sapore opposti. Nonostante notevoli sforzi sperimentali, non sono mai stati osservati quark liberi: diversi studi implicano che essi restino confinati in stati legati, detti adroni. L'intensità di ciascuna interazione è espressa da una costante d'accoppiamento, il parametro adimensionale che ne misura la forza a un dato vertice; nonostante vengano chiamate costanti, queste non sono affatto fisse, ma variano con la scala di energia dei processi considerati. Nel caso della
cromodinamica quantistica (QCD), la costante $\alpha_s$ decresce al crescere
dell'energia scambiata, ovvero al diminuire della distanza fra i quark: alle alte
energie l'accoppiamento diventa debole e i quark si comportano come particelle
quasi libere, regime noto come libertà asintotica, che è ciò che rende possibile
il calcolo perturbativo dei processi di alta energia agli acceleratori. Alle basse
energie, cioè alle distanze dell'ordine delle dimensioni adroniche, $\alpha_s$
cresce fino a rendere impossibile separare due quark: l'energia necessaria per
allontanarli produce nuove coppie quark-antiquark anziché liberarli, regime noto come confinamento. Queste due proprietà sono le stesse manifestazioni ma opposte del andamento di $\alpha_s$ con la scala di energia. Il numero netto di quark, cioè la
differenza fra quark e antiquark, è conservato in tutte le interazioni ed equivale
alla conservazione del numero barionico. Il numero di quark di un dato sapore è invece
conservato soltanto nelle interazioni forte ed elettromagnetica, il che equivale alla
conservazione di $I_3$, della stranezza e degli altri numeri quantici di sapore. Nelle
interazioni deboli il sapore può cambiare, ed è proprio questa proprietà, descritta
dalla matrice CKM, a rendere possibile la violazione di CP discussa nella
Sec. 2.6. Gli adroni si presentano in due configurazioni principali: i mesoni, stati legati quark-antiquark, e i barioni, formati da tre quark (o tre antiquark, nel caso degli
antibarioni); il protone, ad esempio, è composto da due quark $u$ e da un quark $d$.
Negli ultimi anni sono stati osservati anche stati esotici, quali tetraquark e
pentaquark, che non modificano il quadro qui descritto. La ragione per cui
solo certe combinazioni risultano possibili è spiegata dalla QCD attraverso il
concetto di colore: il barione $\Delta^{++}$, costituito da tre quark $u$
con spin paralleli nello stato fondamentale, possiede una funzione d'onda simmetrica
nello scambio di due quark sia nella parte di sapore, che in quella di spin e spaziale. Questo viola il principio di esclusione di Pauli\footnote{Il principio di esclusione di Pauli stabilisce che due fermioni identici non possono occupare lo stesso stato quantico; equivale a
richiedere che la funzione d'onda complessiva di un sistema di fermioni identici sia
antisimmetrica rispetto allo scambio dei due fermioni.}. Per risolvere
questa e altre difficoltà fu introdotto un ulteriore grado di libertà, il colore: ogni quark esiste in tre stati di colore (rosso, verde e blu, con i
corrispondenti anticolori per gli antiquark) e che gli adroni osservabili siano
singoletti di colore, ovvero complessivamente neutri. Il colore è una proprietà
distinta e indipendente dal sapore, ed è la carica dell'interazione forte. Ai fermioni finora descritti si affiancano i bosoni, particelle a spin intero. Quelli
di spin 1, detti bosoni di gauge, sono i mediatori delle interazioni: le particelle che costituiscono la materia interagiscono fra loro scambiandoseli a vicenda, e a ciascuna interazione corrisponde un mediatore caratteristico. L'interazione elettromagnetica è veicolata dal fotone $\gamma$, privo di massa e di carica elettrica, il che gli conferisce un raggio d'azione infinito. L'interazione forte è mediata da otto gluoni, anch'essi privi di massa ma
portatori della carica di colore: ciò fa si che siano autointeragenti (a differenza della bosone $\gamma$ nell'interazione elettromagnetica) a distinguere la QCD dalla QED e a produrre il confinamento dei quark. Pertano, il raggio d'azione resta limitato alle dimensioni adroniche nonostante i mediatori siano privi di massa. L'interazione debole è invece mediata dai bosoni
$W^\pm$ e $Z^0$, massivi rispettivamente attorno a $80$ e $91$~GeV: qui il raggio
d'azione è corto proprio a causa della massa dei mediatori, ed è la stessa ragione per
cui i decadimenti deboli risultano tanto più lenti di quelli forti o elettromagnetici.
I $W^\pm$, portando carica elettrica e sono inoltre gli unici mediatori in grado di
modificare il sapore dei fermioni. L'interazione gravitazionale resta al di fuori di
questo quadro; il suo ipotetico mediatore, il gravitone, non è mai stato osservato, ma
alla scala delle particelle elementari i suoi effetti sono del tutto trascurabili. A completare il contenuto in particelle del Modello Standard vi è infine il bosone di Higgs, unico scalare della teoria, con spin $0$. A differenza dei bosoni di gauge non media alcuna interazione: il suo ruolo è quello di rompere spontaneamente la simmetria elettrodebole, conferendo massa ai bosoni $W^\pm$ e $Z^0$ e, attraverso i termini di Yukawa, ai fermioni stessi. Previsto negli anni Sessanta e osservato ad ATLAS e CMS nel 2012, il meccanismo che porta il suo nome è discusso in dettaglio nella Sec. 2.4. Con questi concetti appena discussi, altre domande che lo SM non riesce a spiegare sono: perché sono raggruppati in tre famiglie? e perché in generazioni di massa?
\begin{table}[htbp]
    \centering
    \begin{tabular*}{\textwidth}{@{\extracolsep{\fill}}cccc@{}}
        \toprule
        \textbf{Generation} & \textbf{Doublet} & $\boldsymbol{Q}$ & \textbf{Mass} \\
        \midrule
        I &
        $\begin{pmatrix}\ \nu_e \\ e \end{pmatrix}_L$ &
        $\begin{matrix} 0 \\ -1 \end{matrix}$ &
        $\begin{matrix} < 0.8~\text{eV} \\ 0.511~\text{MeV} \end{matrix}$ \\
        \midrule
        II &
        $\begin{pmatrix}\ \nu_\mu \\ \mu \end{pmatrix}_L$ &
        $\begin{matrix} 0 \\ -1 \end{matrix}$ &
        $\begin{matrix} < 0.19~\text{MeV} \\ 105.7~\text{MeV} \end{matrix}$ \\
        \midrule
        III &
        $\begin{pmatrix}\ \nu_\tau \\ \tau \end{pmatrix}_L$ &
        $\begin{matrix} 0 \\ -1 \end{matrix}$ &
        $\begin{matrix} < 18.2~\text{MeV} \\ 1777~\text{MeV} \end{matrix}$ \\
        \bottomrule
    \end{tabular*}
    \caption{Lepton doublets of the Standard Model, with electric charge $Q$ in units of $e$.  Neutrino masses are upper limits.}
    \label{tab:sm_leptons}
\end{table}
\begin{table}[htbp]
    \centering
    \begin{tabular*}{\textwidth}{@{\extracolsep{\fill}}ccccc@{}}
        \toprule
        \textbf{Generation} & \textbf{Doublet} & $\boldsymbol{Q}$ & \textbf{Flavour} & \textbf{Mass} \\
        \midrule
        I &
        $\begin{pmatrix} u \\ d \end{pmatrix}_L$ &
        $\begin{matrix} +2/3 \\ -1/3 \end{matrix}$ &
        $\begin{matrix} I_3 = +1/2 \\ I_3 = -1/2 \end{matrix}$ &
        $\begin{matrix} 2.16~\text{MeV} \\ 4.67~\text{MeV} \end{matrix}$ \\
        \midrule
        II &
        $\begin{pmatrix} c \\ s \end{pmatrix}_L$ &
        $\begin{matrix} +2/3 \\ -1/3 \end{matrix}$ &
        $\begin{matrix} C = +1 \\ S = -1 \end{matrix}$ &
        $\begin{matrix} 1.27~\text{GeV} \\ 93.4~\text{MeV} \end{matrix}$ \\
        \midrule
        III &
        $\begin{pmatrix} t \\ b \end{pmatrix}_L$ &
        $\begin{matrix} +2/3 \\ -1/3 \end{matrix}$ &
        $\begin{matrix} T = +1 \\ B = -1 \end{matrix}$ &
        $\begin{matrix} 172.5~\text{GeV} \\ 4.18~\text{GeV} \end{matrix}$ \\
        \bottomrule
    \end{tabular*}
    \caption{Left-handed quark doublets of the Standard Model.
 Nei modelli di bassa energia degli adroni i quark sono
    considerati vestiti a causa della loro interazione con i gluoni che sono i portatori
    della forza forte. Questi quark vestiti sono chiamati quark costituenti, e la massa nella tabella fa riferimento al quark costituente.}
    \label{tab:sm_quarks}
\end{table}



\clearpage

\section{Standard Model}

The Standard Model is a Gauge Theory \(SU(3)_C\ \otimes\ SU(2)_L\ \otimes\ U(1)_Y\) where \(SU(3)_C\) is the group of \(3\times3\) Special Unitary Matrices, associated with an internal colour quantum number. It represents Quantum Chromodynamics (QCD), the strong interaction, where both left-handed and right-handed quarks are subject to the same action of this group; that is, they are affected in a vector-like way, with only the colour component being rotated. The \(SU(2)_L\) is the group of \(2\times2\) Special Unitary Matrices, which acts only on left-handed particles and is therefore a chiral group. Unlike the group \(SU(3)_C\), it does not possess an internal quantum number; consequently, the associated Pauli matrices act by rotating the components of the family doublet. The \(U(1)_Y\), instead, is the Unitary Group of Hypercharge. The "1" indicates that the group has a single generator, corresponding to a phase. This is also a chiral group, since the left-handed hypercharge differs from the right-handed hypercharge. Associated with these groups is the Lagrangian of the Standard Model in the chiral gauge symmetry framework, expressed in the unbroken basis:
\begin{equation}
\mathcal{L}_{SM}\ \supset\ -\frac{1}{2}\mathrm{Tr}\left\{ G_{\mu\nu}(x)G^{\mu\nu}(x) \right\}-\frac{1}{2}\mathrm{Tr}\left\{ B_{\mu\nu}(x)B^{\mu\nu}(x) \right\}-\frac{1}{4}\ C_{\mu\nu}(x)C^{\mu\nu}(x)
\end{equation}
where the first term is a non-Abelian term associated with the group \(SU(3)_C\), describing eight generators:
\begin{equation}
G_\mu(x)=\sum_{A=1}^8 G_\mu^A \hat{T}^A
\end{equation}
where \(A=1,\dots,8\) labels the gluons, expressed through the Gell-Mann matrices. The second term is also a non-Abelian term, associated with the group \(SU(2)_L\) and describing three generators:
\begin{equation}
B_\mu(x)=\sum_{i=1}^3 B_\mu^i \frac{\hat{\sigma}^i}{2}
\end{equation}
where \(\hat{\sigma}^i\) are the Pauli matrices; the last term, instead, is an Abelian term associated with the group \(U(1)_Y\) and has just one generator. Therefore, in the unbroken basis of the SM, we have several vector bosons associated with gauge interactions: 8 vector bosons \(G_\mu\) associated with the strong interaction, and 4 vector bosons \(B_\mu^i\) and \(C_\mu\) associated with the electroweak theory. These gauge bosons are massless, and therefore possess only two degrees of freedom, corresponding to the two transverse polarisations with respect to the direction of motion. However, a model described entirely in terms of massless particles is in clear contradiction with experimental observations, as one can see by simply considering the \(W^{\mp}\) and \(Z^0\) bosons. The resolution of this problem had to await the 1960s, when Peter Higgs and Jeffrey Goldstone solved the problem of the spontaneous breaking of global and local gauge symmetries.

\section{Higgs Potential}
Using the Higgs doublet notation:
\begin{equation}
\phi = \begin{pmatrix} \phi^+ \\ \phi^0 \end{pmatrix}
\end{equation}
introduced to compensate for the quantum numbers of the \(SU(2)_L\) and \(U(1)_Y\) groups within the Yukawa terms, if we include the following terms in the SM Lagrangian:
\begin{equation}
\mathcal{L}_{SM} \supset -\phi^\dagger D_\mu D^\mu \phi +\lambda \phi^\dagger \phi(\phi\phi^\dagger-v^2)
\end{equation}
we can define the Higgs Potential:
\begin{equation}
V^H\left[ |\phi|^2\right] = \int d^4x\ \lambda|\phi|^2(|\phi|^2-v^2) \quad \text{with}\quad v^2=-\frac{\mu^2}{\lambda}>0
\end{equation}
where \(|\phi|^2\) is gauge invariant, and so \(V^H\) is gauge invariant as well. Moreover, the potential follows:
\begin{equation}
V^H\sim|\phi|^4-|\phi|^2
\end{equation}
in full agreement with the renormalisability of the theory, that is, in the classical action there cannot exist operators with engineering dimension greater than the spacetime dimension.

\noindent The Higgs doublet has 4 degrees of freedom (i.e.\ 4 real numbers), that is the real part \(\mathrm{Re}(\phi^+,\phi^0)\) and the imaginary part \(\mathrm{Im}(\phi^+,\phi^0)\). However, it is convenient to switch to polar coordinates:
\begin{equation}
\phi = \begin{pmatrix} \phi^+ \\ \phi^0 \end{pmatrix} = e^{i\ \overrightarrow{\pi (x)}\ \frac{\overrightarrow{\sigma}}{2}} \begin{pmatrix} 0 \\ \frac{v+h(x)}{\sqrt{2}} \end{pmatrix}
\end{equation}
where \(h(x)\) is a fluctuation around the minimum and is a real number; indeed, for \(h(x)=0\), the minimum occurs at \(|\phi|=\pm v/\sqrt{2}\). The field \(\overrightarrow{\pi(x)}\) represents three symmetry rotations of the potential (three real numbers, therefore the number of degrees of freedom is conserved), such that \(V^H\) is invariant under these rotations. In particular, a valley of minima is formed, an infinite set of equivalent minima which, in quantum field theory, correspond to the minima of the Hamiltonian: that is, there are infinitely many vacuum states. These are connected by an \(SU(2)_L\) symmetry operation: the three rotations \(\overrightarrow{\pi(x)}\), which allow transitions from one vacuum state to another without any energy cost. The transition between different vacuum states gives rise to spontaneous symmetry breaking. The term ``breaking'' can however be misleading, because the full theory remains globally invariant under \(SU(2)_L\), and this phenomenon merely hides (rather than removes) the symmetry.

\section{Higgs Mechanism}
The Higgs Mechanism was introduced to explain how a local gauge symmetry can be broken. By fixing the unitary gauge:
\begin{equation}
e^{i\ \overrightarrow{\pi (x)}\ \frac{\overrightarrow{\sigma}}{2}} =1
\end{equation}
and due to locality, it is possible to remove the phase at every point in spacetime, i.e.\ to eliminate the dependence on \(\overrightarrow{\pi(x)}\). This would not have been possible in the case of a global symmetry, since there is no pointwise dependence on \(x\). In this way, the Higgs doublet reduces to (neglecting the \(h(x)\) term for simplicity):
\begin{equation}
\phi = \begin{pmatrix} 0 \\ v/\sqrt{2} \end{pmatrix}
\end{equation}
The kinetic term becomes:
\begin{equation}
-\phi^\dagger D_\mu D^\mu \phi \propto -\frac{1}{2} (0\ \ v) \left\{ \partial_\mu - ig_Y\frac{\hat{Y}}{2}C_\mu(x)-ig_wB^i_\mu(x)\frac{\hat{\sigma}^i}{2} \right\}^2 \begin{pmatrix} 0 \\ v \end{pmatrix}
\end{equation}
and we can neglect \(\partial_\mu\) because \(v=\mathrm{const}\). After a few simple steps, one finds that:
\begin{equation}
W^+=\frac{B^1_\mu-iB^2_\mu}{\sqrt{2}}=\left( W^-_\mu \right)^\dagger \quad, \ Z^0_\mu=- \frac{g_YC_\mu-g_wB^2_\mu}{g} \quad, \ g=\sqrt{g_Y^2+g_w^2}
\end{equation}
with:
\begin{equation}
-\phi^\dagger D_\mu D^\mu \phi \propto\frac{v^2}{2} \left\{ g_w^2 \frac{W_\mu^+W_\mu^-}{2} + g^2\frac{Z^0_\mu Z^0_\mu}{4} \right\}
\end{equation}
where \(\frac{W_\mu^+W_\mu^-}{2}\) is a complex field of the type \(\phi^\dagger \phi\) and \(\frac{Z^0_\mu Z^0_\mu}{4}\) is a real field of the type \(|\phi|^2\). This yields the masses of the weak bosons:
\begin{equation}
m_W^2=v^2\frac{g_w^2}{4} \quad, \ m_Z^2=v^2\frac{g^2}{4}
\end{equation}
where these masses are related through the Weinberg angle:
\begin{equation}
\frac{m_W^2}{m_Z^2}=\cos^2\theta_w\ \quad\ , \ \frac{g_Y^2}{g^2}=\sin^2\theta_w
\end{equation}
from which one experimentally obtains that:
\begin{equation}
m_W\sim 80\ \text{GeV}\ \quad\ , \ m_Z\sim 91\ \text{GeV}
\end{equation}
With this Weinberg angle, it is also possible to relate:
\begin{equation}
\begin{pmatrix} A_\mu \\ Z^0_\mu \end{pmatrix} = \begin{pmatrix} \sin\theta_w & \cos\theta_w \\ \cos\theta_w & -\sin\theta_w \end{pmatrix} \begin{pmatrix} B_\mu^2 \\ C_\mu \end{pmatrix}
\end{equation}
where \(A_\mu\) is the photon, expressed as the combination orthogonal to \(Z^0_\mu\); it remains massless. From the covariant derivative of the quarks, considering only \(B^2_\mu \frac{\hat{\sigma}^3}{2}\), and isolating \(A_\mu\), one finds that the electric charge is:
\begin{equation}
D_\mu q\propto-ig\ \sin\theta_w\ \cos\theta_w\ A_\mu\ \left( \frac{\hat{Y}}{2} +\frac{\hat{\sigma}^3}{2}\frac{1-\gamma^5}{2} \right) = -i\ e\ A_\mu\ \hat{Q}_e
\end{equation}
In the Higgs Mechanism, when a vector boson acquires mass, it gains an additional degree of freedom: the longitudinal polarisation. In our case, three such degrees are required, one for each boson that has acquired mass, and it is natural to ask where these extra degrees of freedom come from. They are the three Goldstone bosons \(\overrightarrow{\pi(x)}\), which, upon reorganisation of the basis, reappear as longitudinal polarisations:
\[\ll\ "W^\mp_\mu\ \text{and}\ Z^0_\mu\ \text{have eaten the Goldstone bosons and got mass.}"\ \gg \]
This is the so-called ``broken basis'' of the Standard Model.

\section{Yukawa Term}
Using the expression of the Higgs doublet in the unitary gauge, let us see what happens if we diagonalise the Yukawa terms:
\begin{equation}
\mathcal{L}_{SM} \supset -\overline{q}_f \left\{ U_{fg}\phi^c\frac{1+\sigma^3}{2}+D_{fg}\phi\frac{1-\sigma^3}{2} \right\}\frac{1+\gamma^5}{2}q_g+h.c.
\end{equation}
where \(\phi^c=i\sigma^2\phi^*\) is the charge conjugate. Collecting all the families into a single large vector \(\overline{q}=(\overline{u},\overline{c},\overline{t},\overline{d},\overline{s},\overline{b})\), explicitly writing the Hermitian conjugate and performing a chiral rotation, we obtain a diagonalised version:
\begin{equation}
\overline{q}=\overline{X}\left\{ \begin{pmatrix} A_U^\dagger & 0 \\ 0 & A_D^\dagger \end{pmatrix} \frac{1+\gamma^5}{2} + \begin{pmatrix} B_U^\dagger & 0 \\ 0 & B_D^\dagger \end{pmatrix} \frac{1-\gamma^5}{2}  \right\}
\end{equation}
\begin{equation}
q=\left\{ \begin{pmatrix} A_U & 0 \\ 0 & A_D \end{pmatrix} \frac{1-\gamma^5}{2} + \begin{pmatrix} B_U & 0 \\ 0 & B_D \end{pmatrix} \frac{1+\gamma^5}{2}  \right\}X
\end{equation}
\begin{equation}
\mathcal{L}_{SM} \supset\frac{v}{\sqrt{2}}\overline{X} \left\{ \begin{pmatrix} A_U^\dagger UB_U & 0 \\ 0 & A_D^\dagger DB_D \end{pmatrix} \frac{1+\gamma^5}{2} + \begin{pmatrix} B_U^\dagger U^\dagger A_U & 0 \\ 0 & B_D^\dagger D^\dagger A_D \end{pmatrix} \frac{1-\gamma^5}{2}  \right\}X
\end{equation}
where the matrices \(A\) and \(B\) are unitary. By the fundamental theorem of algebra, any matrix can always be diagonalised via a biunitary transformation, i.e.\ with real eigenvalues. Therefore, with a suitable choice of \(A\) and \(B\), one can write:
\begin{equation}
\begin{pmatrix} m_u & 0 & 0 \\ 0 & m_c & 0 \\ 0 & 0 & m_t \end{pmatrix} = \frac{v}{\sqrt{2}}\ A^\dagger_U UB_U = \frac{v}{\sqrt{2}} B^\dagger_U U^\dagger A_U
\end{equation}
\begin{equation}
\begin{pmatrix} m_d & 0 & 0 \\ 0 & m_s & 0 \\ 0 & 0 & m_b \end{pmatrix} = \frac{v}{\sqrt{2}}\ A^\dagger_D DB_D = \frac{v}{\sqrt{2}} B^\dagger_D D^\dagger A_D
\end{equation}
So, all the masses in the Standard Model are proportional to \(v\), and the Yukawa terms reduce to a simple sum over the quark masses:
\begin{equation}
\mathcal{L}_{SM} \supset -\overline{q}_f \left\{ U_{fg}\phi^c\frac{1+\sigma^3}{2}+D_{fg}\phi\frac{1-\sigma^3}{2} \right\}\frac{1+\gamma^5}{2}q_g\ + h.c. =\sum_f m_f\overline{X}_fX_f
\end{equation}
By also including \(h(x)\), we can obtain the Higgs interaction vertices, which are proportional to the quark masses, as shown in \cref{fig:hver}:
\begin{figure}[H]
    \centering
    \adjustbox{valign=c}{%
    \(\overline{X}_f\, h(x)\, X_f\)
    }
    \hspace{0.5cm}
    \adjustbox{valign=c}{%
    \includegraphics[width=0.25\linewidth]{TESI/hver.png}
    }
    \caption{Higgs interaction vertices, shown as a Feynman diagram.}
    \label{fig:hver}
\end{figure}
\noindent The change of basis induced by the chiral rotation also affects the kinetic term of the quarks, in particular one term of the weak interactions, giving rise to charged and neutral current interactions:
\begin{equation}
\begin{aligned}
\overline{q}\gamma^\mu D_\mu q ={}& \overline{X}\gamma^\mu\Bigg\{\partial_\mu
  - i(G_\mu,A_\mu,Z^0_\mu) \\[0.4cm]
&\quad - ig_w\frac{\hat{V}_{CKM}W^+_\mu\sigma^+ + (\hat{V}_{CKM})^\dagger W^-_\mu\sigma^-}{\sqrt{2}}\,
  \frac{1-\gamma^5}{2}\Bigg\}X
\end{aligned}
\end{equation}
where \(- i(G_\mu,A_\mu,Z^0_\mu)\) encloses all the neutral current terms, which do not mix the flavour, whereas the second term represents the charged current, which can mix the flavour. The presence of \(\hat{V}_{CKM}\) implies that in weak interactions the charged current interactions mix flavour, since the \(\hat{V}_{CKM}\) matrix gives the probability of mixing between the different quark families (the mixing is larger on the diagonal, especially between the first and second family). The unitarity of \(\hat{V}_{CKM}\), i.e.\ \(\hat{V}_{CKM}(\hat{V}_{CKM})^\dagger=A^\dagger_U A_D A^\dagger_D A_U=\mathbb{I}\), together with its complex phase, gives rise to CP symmetry violation.

\section{CP Violation}
\label{cp_violation}
Weak interactions violate parity (P), charge conjugation (C) and CP symmetry. The latter is the product of the first two and plays a fundamental role in explaining the asymmetry between matter and antimatter in the Universe, although not with sufficient magnitude. Performing a CP transformation on the charged current weak interaction term, we get:
\begin{equation}
\begin{aligned}
\mathcal{L}_{SM} &\supset
  \ccterm{\hat{V}_{CKM}W^+_\mu\sigma^+ + (\hat{V}_{CKM})^\dagger W^-_\mu\sigma^-} \\[0.5cm]
&\overset{CP}{\longrightarrow}\
  \ccterm{(\hat{V}_{CKM})^*W^+_\mu\sigma^+ + (\hat{V}_{CKM})^\dagger W^-_\mu\sigma^-}
\end{aligned}
\end{equation}
The only thing that prevents this term from being invariant is whether:
\begin{equation}
(\hat{V}_{CKM})^*\overset{?}{=}\hat{V}_{CKM}
\end{equation}
If the matrix is real, the term is invariant under CP; otherwise, it is not.
We know that \(\hat{V}_{CKM}\) couples up-type quarks to down-type quarks in charged current weak interactions:
\begin{equation}
\hat{V}_{CKM}=\begin{pmatrix} \hat{V}_{ud} & \hat{V}_{us} & \hat{V}_{ub} \\ \hat{V}_{cd} & \hat{V}_{cs} & \hat{V}_{cb} \\ \hat{V}_{td} & \hat{V}_{ts} & \hat{V}_{tb} \end{pmatrix}
\end{equation}
In the most general case, it contains \(n^2\) complex entries, where \(n\) is the number of families. We may then ask: if it were real, how many angles would it contain? A real \(n\times n\) matrix has \(\#\text{angles}=\frac{n(n-1)}{2}\) and \(\#\text{phases}=n^2-\#\text{angles}=\frac{n(n+1)}{2}\). For \(n=3\), there are three rotation angles and six phases. However, not all of these six phases are physical. By performing a phase transformation of the form:
\begin{equation}
q_u=e^{i\hat{\phi}_u}q_u \quad q_d=e^{i\hat{\phi}_d}q_d
\end{equation}
this affects all the terms in the Lagrangian; however, they are diagonal and therefore remain invariant. The only term that is not invariant is:
\begin{equation}
\hat{V}_{CKM}\longrightarrow e^{-i\hat{\phi}_u}\hat{V}_{CKM}e^{i\hat{\phi}_d}
\end{equation}
This allows us to absorb \(2n-1\) unphysical phases, and therefore the effective number of physical phases is given by:
\begin{equation}
\#\text{Physical phases}= \#\text{Phases} - \#\text{Unphysical phases}= \frac{(n-1)(n-2)}{2}
\end{equation}
Then, for \(n=3\), \(\hat{V}_{CKM}\) is parametrised by 3 rotation angles and 1 complex phase. 
This shows that a physical phase can exist and can't be removed by redefining the fields, but at the same time it doesn't prove that this phase has to be different from zero. Its non-zero value is an experimental fact.
Therefore the Cabibbo--Kobayashi--Maskawa matrix cannot be real, and the presence of this single phase leads to the breaking of CP symmetry in the Standard Model.

\section{Muon in the SM and interaction with matter}







\clearpage

% ==================================================================
%  CAPITOLO 3 - IL RIVELATORE ATLAS
% ==================================================================
\chapter{The ATLAS Detector}

\lettrine[lines=3,loversize=0.1,lraise=0.05]{T}{} he Large Hadron Collider (LHC) is the biggest and most powerful particle accelerator ever built. It’s at CERN, near Geneva, inside a circular tunnel about 27 km around. Almost 100 meters underground, it collides together proton beams (and sometimes heavy ions) at super high energies, letting scientists check out new elementary particles and the fundamental forces that show up from the collisions. There are four main experiments around the ring: ATLAS and CMS (general-purpose detectors), LHCb (for beauty quark physics), and ALICE (looking at heavy ions and quark-gluon plasma). This chapter explains how the various detectors that make up the ATLAS experiment are arranged, going through the reference system adopted by the ATLAS collaboration and focusing the discussion on the muon spectrometer (see Chapter 4), an essential element for analyzing its performance and muon resolution during Run 3. Also, a brief about the ATLAS trigger system will be included in this chapter. 

\section{ATLAS coordinate system}

The reference system used by the ATLAS collaboration places the origin of the axes at the geometric center of the ATLAS detector, that is, at the nominal interaction point (IP). A schematic representation of the axes is shown in \cref{fig:CS}, where the z-axis is oriented along the beam pipe, the y-axis points upwards, and the x-axis runs from the IP to the center of the LHC ring. In the (x,y) plane transverse to the beam, cylindrical coordinates (r, \(\phi\)) are used, where \(\phi\subset [0,2\pi]\) is the azimuthal angle between the x-axis and the particle's transverse momentum. The polar angle \(\theta \subset [0,\pi]\), on the other hand, is the angle between the z-axis and the particle's momentum. Therefore, \(\phi\) helps to understand the particle's position around the z-axis, while theta tells how far it is from it. Although rotations in \(\theta\) around the z-axis are invariant, since ATLAS has cylindrical symmetry, a different coordinate called pseudorapidity \(\eta = -\ ln \left[ tg\left( \theta/2 \right) \right] \) is used; see \cref{tab:eta-theta}. The physical meaning stays the same, but the difference is that for particles coming from the same collision, their cones have different \(\Delta\theta\) but the same \(\Delta\eta\). With this, you can also check out the angular distance \(\Delta R = \sqrt{\Delta\eta^2 + \Delta\phi^2}\).

\begin{figure}[H]
    \centering
    \includegraphics[width=\linewidth]{TESI/CS.jpeg}
    \caption{Schematic representations of ATLAS's coordinate system}
    \label{fig:CS}
\end{figure}
\noindent
\begin{table}[H]
  \centering
  \begin{tabular}{lccccc}
    \toprule
    \(\theta\) & \(90^\circ\) & \(40^\circ\) & \(15^\circ\) & \(5.7^\circ\) & \(0^\circ\)\\
               & \footnotesize(y-axis) & & & & \footnotesize(z-axis) \\
    \midrule
    \(\eta\)   & \(0\) & \(1\) & \(2\) & \(3\) & \(\infty\) \\
    \bottomrule
  \end{tabular}
    \caption{Schematic conversion of the angle \(\theta\) into pseudorapidity \(\eta\)}
    \label{tab:eta-theta}
\end{table}

\section{ATLAS layout during Run 3}

The detector is 25 m high and 44 m long; its overall weight is approximately 7000 t. Based on the coordinate system we just introduced, let's try to understand how the various ATLAS detectors are distributed. A quick overview of the Run 3 configuration is shown in \Cref{fig:ATLASexp}. Moving outwards from the beam line, within the 2 T xial magnetic field of the solenoid, we have: the Inner Detector system (ID) that allow us to reconstruct the tracks of the charged particles within \(|\eta|<2.5\).this is composed by the High-Granularity silicon pixel detector that covers the vertex region providing four measurements per track, where generally the first hit start in the Insertable B-Layer (IBL) \cite{ATLAScollab2, ATLASibl}, outside them we have the Semiconductor Tracker (SCT), a silicon microstrip tracker that provides eight measurements per track. Moreover, these silicon detectors are integrated with Transition Radiation Trackers (TRTs) that provide information on the electronic component within \(|\eta|=2.0\). Outside the solenoid, the calorimeter system gives full solid-angle coverage, going up to \(|\eta|<4.9\). The electromagnetic calorimeters sit in the \(|\eta|<3.2\) region, especially in the barrel and endcap, and they’re high-granularity, made of lead or Liquid Argon (LAr). As for the hadronic calorimeters, three are in the \(|\eta|<1.7\) area (barrel), made of alternating steel and scintillator layers, and two are in the \(1.5<|\eta|<3.2\) region (endcap), made of alternating copper and LAr layers. To cover the rest of the solid angle, the forward direction has smaller calorimeters in the same setup. The Muon Spectrometer (MS) comprises separate trigger and high-precision tracking chambers in order to measure the bending of muon tracks in the field produced by the superconducting air-core toroidal magnets, whose field integral lies between 2.0 Tm and 6.0 Tm over most of the detector. Then, we have three stations of precision chambers, where the innermost one of the end-caps, within \(|\eta|>1.3\), has been replaced by the New Small Wheel (NSW) detectors for Run 3. The other two stations, instead, each consist of layers of Monitored Drift Tubes (MDTs), which allow us to reconstruct tracks in the barrel. Installed during the Long Shutdown~2 in preparation for Run~3, the NSW is built with Micromegas (MM) and Thin Gap Chambers (sTGC), and supersedes the former Small Wheel instrumentation of MDTs, Cathode Strip Chambers (CSCs), and a single station of TGCs. The muon trigger system is provided, over the range $|\eta| < 2.4$, by Resistive Plate Chambers (RPCs) in the barrel ($|\eta| < 1.0$), while Thin Gap Chambers (TGCs) cover the endcap regions ($|\eta| > 1.0$), outside the toroid magnets. An additional contribution comes from the Micromegas (MM) and small-strip TGC detectors of the New Small Wheels, located between the cryostats of the endcap calorimeters and the endcap toroids. Interesting events are selected by a level-1 trigger system built in custom hardware, and are then filtered further by software algorithms running in the high-level trigger~\cite{ATLAScollab3}. Out of the 40~MHz bunch-crossing rate, the level-1 system retains events at a rate of at most 100~kHz; the high-level trigger applies additional selections so that data are written to disk at approximately 3~kHz. The configuration of the ATLAS detector during Run~3 is shown in \Cref{fig:ATLASexp}.
\\[2cm]
\begin{figure}[H]
    \centering
    \includegraphics[width=1\textwidth]{TESI/ATLASexp.png}
    \caption{Section of the Run 3 configuration of the ATLAS detector showing the locations of the larger detector sub-systems.}
    \label{fig:ATLASexp}
\end{figure}

\section{The Trigger System}
The ATLAS Trigger and Data Acquisition (TDAQ) system is responsible for selecting interesting events for analysis and reading them for offline processing, such as energetic leptons, photons, hadronic jets, tau leptons, or large missing energy. The TDAQ is organized into two selection levels: Level-1, made with custom electronics, and the High-Level Trigger, which is a software-based system running on commercial computers. These two levels share the Data Acquisition (DAQ) system, which carries data from the sub-detector electronics to offline processing following the trigger decisions. The L1 works on detailed information coming from the calorimeters and the muon spectrometer. Its maximum rate of acceptance sustainable by the readout systems is 100 kHz, while all processing has to be completed within the latency allowed by the detector electronics, which is \(2.5\ \mu s\) per event. The data acquisition system has also been upgraded so it can read the whole detector during Run 3. The ReadOut System needs to handle 100 kHz of events accepted by Level-1, plus an extra \(20\%\) margin. In Run 3, the load it needs to manage increases for two reasons: HLT asks for data more often, and the average event size goes up by \(30\%\), reaching \(2.1\ MB\) at \(\langle\mu\rangle \approx 60\). Further down, the system has to write an average flow of 3 kHz to mass storage while staying flexible enough to handle rate spikes. The max available bandwidth is 8 GB/s: in the example given, with \(\langle\mu\rangle \approx 60\) and a max physics event rate on disk of \(3.3\ kHz\), it needs \(6.9\), almost double compared to Run 2. In the end, to record raw data in continuous data acquisition at 3 kHz, at Point 1 there is at least \(1.4\ PB\) of effective space available: a capacity that ensures 24 hours of autonomy in case of downtime of CERN's mass storage services and accounts for both the time needed to transfer and then delete the files, as well as the margin required for the file system to maintain expected performance.

\clearpage

% ==================================================================
%  CAPITOLO 4 - MUON SPECTROMETER
% ==================================================================
\chapter{The ATLAS Muon Spectrometer}
\lettrine[lines=3,loversize=0.1,lraise=0.05]{H}{} igh momentum muons in the final state are among the most promising and robust signatures of the physics accessible at the Large Hadron Collider (LHC). Being the only charged particles that traverse the calorimeters almost undisturbed, they can be identified with little ambiguity even in the dense environment of \(pp\) collisions. To exploit this potential, the ATLAS Collaboration has designed a high-resolution Muon Spectrometer, able to trigger on muons and to measure their momentum in stand-alone mode over a wide range of transverse momentum, pseudorapidity and azimuthal angle. This chapter gives a brief description of the main components of the spectrometer, of the role each of them plays and of their arrangement within the system. Five detector technologies are employed for triggering and precision tracking: Resistive Plate Chambers, Thin Gap Chambers, Monitored Drift Tubes, Micromegas and small-strip TGCs, and each is presented in the following sections.


\section{Muon System}
The Muon Spectrometer (MS)~\cite{ATLASdesc} constitutes the outermost and largest part of the ATLAS experiment. Its role is to detect the charged particles that emerge from the barrel and endcap calorimeters, allowing us to measure their momentum up to \(|\eta| = 2.7\). The measurement relies on the deflection of muon tracks inside large superconducting air-core toroid magnets, complemented with two sets of chambers, one dedicated to triggering and one to high-precision tracking. Magnetic bending is provided by the large barrel toroid for \(|\eta| < 1.4\) and by two smaller endcap toroids, inserted at either end of the barrel magnet, for \(1.6 < |\eta| < 2.7\); in the transition region between them, muons are deflected by the combined barrel and endcap fields. This arrangement of magnets keeps the field orthogonal to the muon tracks, ensuring a smaller loss of resolution induced by multiple scattering. The MS joins the trigger system up to \(|\eta| = 2.4\). The MS is organised into a barrel and two endcaps. Three concentric stations of roughly cylindrical shape make up the barrel and are labelled: inner (BI), middle (BM) and outer (BO). Each endcap is, instead, built from three disc-shaped stations, known as wheels and labelled inner (EI), middle (EM) and outer (EO), to which an Extended Endcap ring (EE) is added at a radius larger than that of the endcap toroid cryostat. The three wheels occupy rather different positions along the beam line. The EI wheels (New Small Wheels) are housed within the volume enclosed by the barrel toroid, in the gap that separates the calorimeters from the endcap toroid cryostats. The EM wheels (Big Wheels) lie immediately downstream of those cryostats, while the EO wheels, with the same size of the EM ones, are carried by metal support structures anchored to the end-walls of the ATLAS cavern. The role of the EE rings is to add a third measurement point at intermediate pseudorapidity, \(1.05 < |\eta| < 1.3\), for the tracks that escape the coverage of the EO wheel; the rings are accordingly placed between the EI and EM wheels. A special case is that of the EIL4 stations, which cover the interval \(1.0 < |\eta| < 1.25\): instead of being carried by the EI wheels, they occupy a fixed position among the coils of the barrel toroid. As far as the barrel instrumentation is concerned, the bending coordinate is measured in all three stations by multilayer Monitored Drift Tube (MDT) chambers; the BM and BO stations carry, in addition, Resistive Plate Chambers (RPCs), from which both the trigger decision and the coordinate along the azimuthal direction are obtained. Only MDT chambers are installed in the outer wheels, while the middle wheels combine MDTs, for precision tracking in the bending plane, with TGCs, used both for triggering and for the azimuthal measurement. The inner wheels are located between the calorimeters and the endcap toroid cryostats, within the volume of the barrel toroid. They have been entirely replaced by the NSWs, which occupy the same location and cover the same polar range, \(1.3 < |\eta| < 2.7\). The NSWs use two chamber technologies: small-strip TGC (sTGCs) and micro-mesh gaseous structure (Micromegas) chambers. Both are fast enough to contribute to the Level-1 decision, and both are designed for precision tracking along the bending direction as well as for improved resolution in the azimuthal coordinate. The resulting gain in trigger \(p_T\) resolution and in background rejection makes it possible to retain a low muon \(p_T\) threshold at a sustainable Level-1 rate, thereby preserving the acceptance for many physics processes of interest. The layout of the MS is shown in \cref{fig:MS}, and its main parameters are summarised in \cref{tab:MS}.

\begin{figure}[H]
    \centering
    \includegraphics[width=1\textwidth]{TESI/MS_nuovo.png}
    \caption{Section of the ATLAS muon system, built around an air-core toroid magnet system, which measures the momentum of the charged particles leaving the barrel and endcap calorimeters within \(|\eta| < 2.7\); the components upgraded for Run 3, in particular the NSWs, are shown.}
    \label{fig:MS}
\end{figure}

\begin{table}[H]
    \centering
    \includegraphics[width=0.8\textwidth]{TESI/tab_04.pdf}
    \caption{Main parameters of the muon spectrometer}
    \label{tab:MS}
\end{table}

\section{Monitored Drift Tube (MDT)}
Monitored Drift Tube (MDT) chambers constitute a fundamental part of the precision tracking system of the ATLAS muon spectrometer, where they measure the trajectories of the charged particles produced in \(pp\) collisions over the pseudorapidity range \(|\eta| < 2.7\). They provide an average spatial resolution of about \(80\ \mu m\) for a single tube, which improves to roughly \(35\ \mu m\) once the measurements of a full chamber are combined. The chambers sustain hit rates as high as 300 kHz per tube with no appreciable loss of either spatial resolution or detection efficiency \cite{ATLASdesc, MDTTDCsim}.  Each chamber is equipped with up to 18 front-end boards, known as mezzanine cards. A mezzanine hosts three custom 8-channel Amplifier/Shaper/Discriminator (ASD) chips and one 24-channel Time-to-Digital Converter (TDC), the ATLAS Muon TDC (AMT), currently in its third version (AMT-3). The signals of all mezzanines belonging to a chamber are collected by an on-chamber data acquisition board, the Chamber Service Module (CSM), which formats and serialises them and transmits them over an optical link to the off-detector MDT Readout Driver (MROD) \cite{MDTTDCsim}. The read-out system was originally dimensioned for an instantaneous luminosity of \(10^{34}\ cm^{-2}s^{-1}\), while Run 3 has already exceeded \(2 \times 10^{34}\ cm^{-2}s^{-1}\) at a Level-1 accept rate of 100 kHz; dedicated simulations indicate that the legacy AMT-3 can cope with these conditions, the hit loss in the busiest chamber remaining below \(5\ \%\) \cite{MDTTDCsim}. For the HL-LHC, where the luminosity may reach \(7 \times 10^{34}~cm^{-2}s^{-1}\) and the Level-0 system is required to accept events at 1~MHz within a \(10\ \mu s\) latency, the front-end electronics will be replaced during Long Shutdown~3 and the new TDCs will operate without trigger mode.

\section{Resistive Plate Chamber (RPC)}
Muon identification in the barrel of ATLAS relies on Resistive Plate Chambers (RPCs) \cite{ATLASRPCRun3}, which deliver both a stand-alone trigger decision and a measurement of two coordinates. Each chamber consists of two independently operating detector layers, and the chambers are distributed over three concentric cylindrical shells located between roughly 7 and 10 m from the beam axis, immersed in a toroidal field of about 0.5 T. Every gas volume is built from bakelite electrodes enclosing a 2 mm gap, continuously flushed with a ternary mixture of \(C_2H_2F_4\), \(\text{iso - }C_4H_{10}\) and \(SF_6\) in proportions of \(94.7\%\), \(5.0\%\) and \(0.3\%\). Operation takes place in avalanche mode: the applied voltage of 9.6 kV corresponds to a field of 4.8 kV/mm across the gap, and it is rescaled automatically to compensate for deviations of temperature and pressure from the reference values of \(24\ ^\circ C\) and 970 mbar \cite{GA}. Signals are collected by readout panels whose strips, between 2.3 and 3.5 cm wide, are oriented so as to provide the two orthogonal coordinates \(\eta\) and \(\phi\). Altogether the system comprises 3720 gas volumes, spans 3650 \(m^2\) of active surface and is read out through approximately 385000 channels. A first upgrade of this layout has already been carried out: 16 RPC triplets (BIS78) were mounted in the transition region between barrel and end-cap, where the limited instrumentation in the interval \(1.0 < \eta < 1.3\) was responsible for a sizeable rate of fake triggers. In these new units, the bakelite electrodes are separated by a narrower 1 mm gap, and the copper readout strips are coupled to low-noise front-end electronics with high amplification. Further work is foreseen for the HL-LHC phase. The innermost barrel layer, at present equipped with drift tubes alone, will be complemented by a triplet of RPC layers, thereby extending the acceptance and adding redundancy to the muon trigger.

\section{Thin Gap Chamber}
As discussed in section 4.1, the TGCs are located in the middle wheels, providing us with Level-1 trigger information and measuring the azimuthal coordinate of muons in the endcaps of the ATLAS MS \cite{TGC1}. TGCs are multi-wire proportional chambers with the characteristic that the wire-to-cathode distance  (1.4 mm) is smaller than the wire-to-wire distance (1.8 mm). The cell structure allows us to measure in a quasi-saturated mode due to a highly quenching gas mixture of \(CO_2\) and \(\text{n -} C_5H_{12}\) (n-pentane), i.e. with a gas gain of \(3\ \cdot\ 10^5\) that does not allow to make full use of its independence from the primary ionisation. Moreover, the highly quenching gas prevents the occurrence of streamers in all operating conditions. Instead, the presence of a high electric field around the TGC wires and, thanks also to the small wire-to-wire distance, leads to an incredible time resolution for the majority of the tracks. Including the variation of the propagation time on wires and strips, signals arrive with \(99\%\) probability inside a time window of 25 ns. More information on the fundamental parameters of TGC construction is shown in the TGC parameter book \cite{TGC}.

\section{New Small Wheel}
The NSW detectors allow us to reconstruct muon tracks with high precision and provide information for the L1 trigger \cite{ATLASdesc, NSW}. As discussed in Section 4.1, they use two peculiar chambers: small-strip TGC (sTGC) and micro-mesh gaseous structure (Micromegas) detectors. Both can be used for the trigger given the high ability of the sTGCs to identify single Bunch Crossings and the high speed of the MM thanks to their small drift gap. Also, they are great for tracking because of the high resolution (\(100\ \mu m\) in the bending direction. These two different detectors reach high spatial resolution with different mechanisms: the MM exploits a very thin strip pitch (\(<0.5\ mm\)), while the sTGCs, having a larger pitch (\(3.2\ mm\)), rely on the charge sharing between the strips. This precision needs to preserve the high muon momentum resolution even in a highly background enviroment as LHC. Moreover, improve a better resolution for the azimuthal coordinate with respect to the EI wheels: better resolution at trigger level and offline reconstruction, in view of Run 3. All this leads to measuring the transverse muon momentum in the endcaps that, combining the track segments in the NSW with those from the ID, must have a better precision of \(15\%\) for a 1 TeV muon all over \(|\eta|<2.7\). To do this, NSW must have a resolution of \(50\ \mu m\) in the bending plane to reconstruct these track segments, independently of the background conditions. Since the NSW has 16 sectors, this leads to a resolution of \(150\text{-}175\ \mu m\) for each strip plane. In fact, from a geometric point of view, each wheel maintains the spectrometer's eightfold symmetry and is divided into eight Large sectors and eight Small sectors (see \cref{fig:NSWa}, instead, the \cref{fig:NSWb} shows the shielding disk JD). In the end, the layout reflects the division of tasks between the two technologies: the sTGCs are on the outside because that's where their angular precision is most useful for the trigger, while the Micromegas occupy the central volume because they're responsible for precise tracking measurement in the curvature coordinate.

\begin{figure}[htbp]
    \centering
    \begin{subfigure}[b]{0.48\textwidth}
        \centering
        \includegraphics[width=\textwidth]{TESI/NSWa.png}
        \caption{}
        \label{fig:NSWa}
    \end{subfigure}
    \hfill
    \begin{subfigure}[b]{0.48\textwidth}
        \centering
        \includegraphics[width=\textwidth]{TESI/NSWjd.png}
        \caption{}
        \label{fig:NSWb}
    \end{subfigure}
    \caption{New Small Wheels: \textit{(a)} shows the structure: the wheel on the left is seen from the cavern end-wall, with the A-frame and shielding hub visible, along with the eight Large sectors and the eight spokes supporting them; the wheel on the right is seen from the IP, and since the newJD shielding disc is not shown, the eight Small sectors, aligned with the barrel toroid coils, are visible, along with the spokes supporting them. \textit{(b)} Support structures for the NSWs, including the newJD discs. }
    \label{fig:nsw-layout}
\end{figure}

\subsection{small-strip TGC (sTGC)}
The sTGCs \cite{ATLASdesc, NSW} were created to do two things at once: provide tracking fast enough to enter the level 1 hardware trigger and, at the same time, a position measurement precise enough to also be used in offline muon reconstruction. Classic TGCs were only used for the trigger; the new thing here is that the same detector has to do both jobs, in an environment with higher background. The technological base remains that of the TGCs: a \(2.8\ mm\) gap bordered by two FR4 cathodes coated with a resistive mixture of graphite and resin, filled with a mixture of \(55\%\ CO_2\) and \(45\%\) n-pentane; in the center of the gap, the anode plane consists of \(50\ \mu m\) gold-plated tungsten wires spaced \(1.8\ mm\) apart. The changes, on the other hand, only concern what was needed for the two new goals. The wires are now parallel to the radial axis of the chamber, whereas in the TGC they are stretched in the azimuthal direction; the readout strips are much denser (\(2.7\ mm\) wide, \(3.2\ mm\) pitch), with a geometry experimentally optimized to limit capacitive coupling between channels and nonlinear corrections, while keeping the number of readout channels as low as possible; finally, the signal layer has been moved under a thin pre-preg, just below the graphite instead of on the back of the FR4, increasing the capacitance to maintain the same transparency to fast signals. All of these choices result in an intrinsic resolution of about \(45\ \mu m\) for particles hitting perpendicularly \cite{NSW2}. To deal with background, they adjusted the resistivity of the cathode coating, lowering it to \(150\ k\Omega/\Box\)\footnote{Not to be confused with resistance, surface resistivity is the resistance between two opposite sides of a square of a thin-film resistor, and it doesn't depend on the size of the square; it's, anyway, measured in \(\Omega\).} in the chambers closest to the interaction point and to \(200\ k\Omega/\Box\) in the others, so the accumulated charge is quickly dispersed. Together with reducing the decoupling resistors, this keeps the chambers efficient for minimally ionizing particles even at rates up to \(20\ kHz/cm^2\) over large areas.

\subsection{micro-mesh gaseous structure (Micromegas)}

Micromegas technology \cite{ATLASdesc} allows the construction of thin planar gas detectors where the traditional high-voltage wire planes are replaced by a metal micro-mesh. The detector consists of a planar drift electrode, a gas volume of a few millimeters that serves as both the conversion and drift region, and the micro-mesh, placed at a distance of \(120\text{-}130\ \mu m\) from the readout electrode, which defines the amplification region. The potentials are chosen so that the electric field is a few hundred V/cm in the drift region and between 40 and 50 kV/cm in the amplification region: the first is adjusted to bring the drift velocity near its local maximum and plateau, thus minimizing longitudinal diffusion and ion backflow; the second is tuned according to the gas mixture to achieve a gain close to 10000. A charged particle passing through the drift volume ionizes the gas; the freed electrons drift toward the mesh, which ends up being transparent to more than \(95\%\) of them precisely because the amplification field is 50 to 100 times stronger than the drift field. The avalanche develops in the thin gap above the readout electrode. The two timescales involved are very different: drift is a relatively slow process, depending on the gas, drift distance, and electric field, and in ATLAS setups it can take up to about 100 ns, while amplification happens in a fraction of a nanosecond, giving a very fast pulse of electrons on the strip. The ions produced in the avalanche move in the opposite direction, towards the mesh, but since most of them are generated in the last stage of multiplication, that is, very close to the readout strip, they are quickly cleared out: it’s this feature, more than anything else, that makes Micromegas suitable for operating at very high particle fluxes. The main challenge in the ATLAS project is discharges in large-area detectors. These discharges kick in when the total number of electrons in an avalanche hits a few tens of millions, reaching the Raether limit, and they can damage both the detector and the readout electronics. The tricky part is that the efficiency needed for minimum ionizing muons requires a gain of around 10,000, so you only need ionization processes that release more than a thousand electrons over distances similar to the lateral spread of an avalanche for the risk to actually come into play. The protection system developed involves adding a layer of resistive strips on top of the readout strips, separated by a thin insulator: the readout electrode is no longer directly exposed to the charge created in the amplification region, but is capacitively coupled to it. You lose a fraction of the signal amplitude, but the chamber can be operated at a higher gain, and the intensity of discharges is reduced by about three orders of magnitude.

\begin{figure}[H]
    \centering
    \includegraphics[width=0.8\linewidth]{TESI/MM.png}
    \caption{Schematic view of an ATLAS NSW Micromegas detector. The micro-mesh is integrated into the drift panel and stretched over the grid of support pillars on the readout panel when the two panels are pressed together. We have a conversion and drift gap of 5 mm, an amplification gap of \(128\ \mu m\), a drift electrode at \(\text{-}300\ V\), mesh at ground, and resistive strips at \(\text{+}505\ V\), with a readout pitch of \(425\ \mu m\). Incoming muons are drawn as red lines; the large red markers indicate where the primary ionisation clusters are produced, the yellow shapes trace the spread of the released electrons as they drift, and the small red markers within the amplification gap correspond to the successive multiplication steps of the avalanche.}
    \label{fig:MM}
\end{figure}

\clearpage

%===================================================================
% CAPITOLO 5 - RICOSTRUZIONE E CLASSIFICAZIONE MUONI
%===================================================================
\chapter{Muon Reconstruction \& Classification}


\lettrine[lines=3,loversize=0.1,lraise=0.05]{I}{} n ATLAS, muon reconstruction is based on information from several detector systems, including the Inner Detector (ID), Muon Spectrometer (MS), and calorimeters \cite{Journ}. Muons are highly penetrating particles (i.e. minimum-ionising particles), identified using tracks in the MS. In order to ensure high efficiency and precision, two main strategies are adopted: the first, discussed in Section 5.1, uses stand-alone track reconstruction in the MS; the second, described in Section 5.2, integrates all the information obtained from the whole ATLAS detector. Once the reconstruction is completed, only high-quality muons are selected to preform the analysis. In particular, they are chosen using several requirements concerning the number of hits in the different ID subdetectors and the number of hits in each MS stations, the quality, and several other properties of the muon track. In the following, several sets of criteria are discussed.

\section{Stand-alone tracks reconstruction}
The stand-alone track reconstruction in the MS is a fundamental process to obtain the characteristics and trajectories about the muon that pass through the detector. This process occurs in several steps that involve combinations of geometrical and physical methods to guarantee a high quality and reliable muon track reconstruction. It starts with the short straight-line local track segments derived from "hits" in an individual MS station (there are typically 3 MS stations in most of the spectrometer). Each segment represents a part of the muon's trajectory, and their identifications are performed using the Hough transform. Once these segments are identified in the different MS stations, they are combined to form preliminary candidate muon tracks using a loose pointing constraint based on the interaction point (IP) and the parabolic trajectory that constitutes a first-order approximation to the muon bending in the magnetic field. Then, Trigger chambers provide a second coordinate (along another axis) in order to obtain a 3D description of the trajectory. Afterwards, a \text{$\chi^{2}$} global fit is performed to the muon's trajectory under the magnetic field. This is a statistical technique used for minimizing the difference between the observed trajectory (the hits on MS) and the expected trajectory, which can have errors due to misalignments in the different detector chambers as well as the several interactions with the detector that can alter the real muon path. From the fit, two operations starts, that is the removal of outlier hits and the addition of the hits along the trajectory that were not assigned to the original track candidate in order to improve the reconstruction. Last but not least, if some ambiguities are present between tracks, for example, two tracks that share the same hits fraction, the lower-quality one is rejected. One exception to this rule is made in the case of tracks that are identical in two stations but share no hits in a third station. This helps to guarantee a high efficiency in the light-flavour muons system reconstruction tracks originating from a second vertex due to the decay of pions or kaons, where the detector can mistakenly interpret it as a true signal from a prompt muon track. The final set of tracks is re-fitted with a loose IP constraint, taking into account the energy loss in the calorimeters, and back-extrapolated to the beam line. The \( p_T \) of the extrapolated track is then expressed at the IP.

\section{Global muon reconstruction}
The global muon reconstruction is based on a system that involves all the detectors in ATLAS experiment, as the ID, the MS and calorimeter systems. Also here, one wants to achive the same result, i.e. trajectory description and properties, but improved by all the info in the detector. There are five strategies to obtain the reconstruction, each producing different muon types. Combined Muons (CB) are reconstructed by comparing tracks that come from MS with those coming from ID. Then a fit is made using the MS and ID hits, taking into account the energy loss in the calorimeters, which can influence the muon path. For \text{$|\eta| > 2.5$}, muon tracks in MS can be associated with partial information from Pixel and SCT detectors, creating a subcategory of CB muons called Forward Muons associated withSilicon Forward muons (SiF). Inside-Out muons, where the tracks are reconstructed starting from ID and then extrapolated to MS, improving the efficiency in regions where MS is limited and for muons that have low \( p_T \). If an MS track cannot be matched to an ID track, its parameters are extrapolated to the beamline and used to define an Extrapolated Muon (ME). This expands the coverage of the MS to \text{$|\eta| = 2.7$}, improving the muon identification in regions where ID failed. Segment-Tagged Muons (ST) are reconstructed starting from an ID track, which is identified as a muon because it is associated with at least one segment in the MS. Unlike the previously mentioned categories, ID and MS information are not combined here; only ID information is used, while MS is used as support to select the tracks. ST muons are used to recover some efficiency in regions where the MS has limited acceptance, and it's not possible to reconstruct even a partial track. Calorimeter-Tagged Muons (CT), as the name suggests, are muons reconstructed starting from ID tracks, which are associated with energy loss in the calorimeters compatible with minimum ionising particles. Also in this case, only the Inner Detector (ID) is used to reconstruct the track, while the calorimeter serves as support to recover efficiency in regions where the Muon Spectrometer (MS) has no acceptance.

\section{Data Samples}

The analysis presented in this chapter is based on Monte Carlo simulated samples: the studied observable requires a comparison between reconstructed momentum and the truth momentum of the muons. All samples used come from the production of \texttt{mc23e} (Athena MC release 2024) for \(pp\) collisions at \(\sqrt{s}=13.6\ TeV\), generated to reproduce the data-taking conditions of the LHC during Run 3 with the pile-up and detector conditions relating to 2024. The response of the detector is simulated with the toolkit \textsc{GEANT4} \cite{geant4} and the events are reconstructed with the same software chain applied to the real data. Two physics process categories are implemented in this analysis. The first one is the Drell-Yan production of the Z bosons with decay into two muons, $Z\rightarrow\mu^+\mu^-$ (DSID 601190), generated with \textsc{Powheg Box} at the Next-to-Leading-Order (NLO) accuracy in QCD and \textsc{Pythia 8} for matrix element (ME) calculations, for the modelling of the parton shower, hadronisation, and the underlying event, with the parameter sets AZNLO; the effect of QED final-state radiation was simulated with \textsc{Photos++}. This sample provides us with high isolated muon statistics with high quality in the region of the Z peak (i.e. \(p_{\mathrm{T}}\simeq45\ GeV\)), that is, the region in which the ATLAS collaboration determines the muon calibration constants \cite{ATLAS-MCP-2023}. The second one is made up of a heavy resonance of \(Z'\rightarrow\mu^+\mu^-\) at zero width (DSID 801862-801866), generated with \textsc{Pythia 8} using a parton density function (PDF) set at the Leadin-Order (LO) level of accuracy and the A14 set of tuned generator parameters were used for five mass values: 500, 1000, 3000, 5000 and 8000 GeV as shown in \cref{fig:peaks} and \cref{tab:peaks}. 
\begin{figure}[htpb]
    \centering
    \includegraphics[width=\textwidth]{TESI/plot_pt_Z_Zprime_minimal.pdf}
    \caption{Distribution in $p_{\mathrm{T}}$ of the reconstructed muons through the Z sample and the five Z' samples at different masses. Each histogram is normalized to its own integral. The value in each bin is therefore the fraction of muons in that sample that fall into that bin of $p_{\mathrm{T}}$: $N_{bin}/N_{tot}$. It serves to compare the shape of the six distributions regardless of the number of muons in each sample. Each curve has a peak at about half the mass of the resonance, that is, the kinematics of a two-body decay at rest, $p^{\mu}_{\mathrm{T}}\simeq M/2$}
    \label{fig:peaks}
\end{figure}
\noindent
\begin{table}[htbp]
    \centering
    \begin{tabularx}{\textwidth}{X|X|X|X|X}
        \toprule
        Sample & Resonance [GeV] & $p_{\mathrm{T}}^{\mathrm{expected}}$ [GeV] & $p_{\mathrm{T}}^{\mathrm{observed}}$ [GeV] & $\Delta p_{\mathrm{T}}/p_{\mathrm{T}}$ \\
        \midrule
        $Z \to \mu^+\mu^-$ & 91.2 & $\simeq 45$   & $\simeq 42$   & $7\%$  \\
        $Z'$ 500 GeV       & 500  & $250$         & $\simeq 244$  & $2\%$  \\
        $Z'$ 1000 GeV      & 1000 & $500$         & $\simeq 455$  & $10\%$ \\
        $Z'$ 3000 GeV      & 3000 & $1500$        & $\simeq 1284$ & $17\%$ \\
        $Z'$ 5000 GeV      & 5000 & $2500$        & $\simeq 2156$ & $16\%$ \\
        $Z'$ 8000 GeV      & 8000 & $4000$        & $\simeq 3263$ & $23\%$ \\
        \bottomrule
    \end{tabularx}
    \caption{Comparison between the expected transverse momentum of the muons, $M/2$, and the position of the peak actually observed in each simulated sample: the Standard Model $Z \to \mu^+\mu^-$ and the five zero-width $Z' \to \mu^+\mu^-$ resonances, whose spectra are shown in \cref{fig:peaks}. Since the resonances are generated at zero width, each sample provides two muons with fixed momentum $M/2$ in the rest frame of the resonance, and the corresponding peak is expected at that value in the transverse plane. The last column gives the relative shift $\Delta p_{\mathrm{T}}/p_{\mathrm{T}} = (p_{\mathrm{T}}^{\mathrm{expected}} - p_{\mathrm{T}}^{\mathrm{obs}})/p_{\mathrm{T}}^{\mathrm{obs}}$. The observed peak lies systematically below the expectation, and the shift grows with the resonance mass, from $2\%$ at 500~GeV up to $23\%$ at 8~TeV.}
    \label{tab:peaks}
\end{table}
\noindent Since the resonance width is zero, each sample provides us with two muons with fixed momentum to \(m_{Z'}/2\) in the resonance reference frame at rest; in the lab reference frame, the transverse muon momentum distribution presents the usually jacobian peak at that value, allowing us to have a full popolate spectrum since the \(Z\rightarrow\mu^+\mu^-\) process have not enough statistics. The same strategy is used in \cite{ATLAS-MCP-2023}, where narrow resonances with a mass over 500 GeV are used to check the validity of the corrections outside the kinematic range in which they were derived. Each process was built in two different configurations of the detector geometry. The first, in the following denoted as \emph{ideal} geometry, assumes a perfect alignment of both the ID and MS: this is the standard condition for the ATLAS simulations, in which no residual misalignment effect appears. The second, denoted as \texttt{MSmisalign}, the chamber positions in the intermediate layer of the MS were deliberately shifted by amounts compatible with the residual precision achieved by the alignment procedures \cite{ATLAS-MCP-2023}, reproducing a realistic chamber shift, which is a distorted geometry that shows their misalignment. The comparison between the two configurations, given the same generated process and applied selection, is the method used in Section 6.5 to isolate the contribution of the residual misalignment. From a technical view, the simulated data samples used in this analysis are \emph{ntuples} \texttt{ANALYSIS.root} (production with \texttt{MCP\_TESTNTUP}), given by files xAOD through the package \texttt{MuonxAODAnalysis}. These ntuples contain, for each selected muon, kinematic characteristics reconstructed separately from the track in the ID, in the MS, and from the combined track (CB). The files are read remotely via the XRootD protocol, from a "txt" file with their paths. The sample of $Z\rightarrow\mu^+\mu^-$, ensuring enough statistics in all the pseudorapidity range, is subdivided into 68 files in perfect alignment configuration and 26 files in misaligned configuration. Instead, for each \(Z'\), one file is enough, for a total of \(\simeq 250 \) million events. 

\section{Prompt muon selection}

Muons can be divided into two classes. \emph{Prompt muons} are produced directly at the primary vertex and are more isolated in the ID and/or the calorimeters. \emph{Non-prompt muons} are instead classified into light flavour muons, coming from the decay of light hadrons such as pions and kaons produced in a gluon jet, and heavy flavour muons, coming from the decay of B- and C-hadrons. As previously discussed, a more accurate measurement will be obtained considering only prompt muons. A prompt muon is produced directly at the primary vertex, leaving a clean track that can be easily reconstructed; it comes principally from the \(W\), \(Z\) and \(Higgs\) boson decays, or from other electroweak processes. In ATLAS Truth Classification, which is a tool that combines the type and origin truth of the particle to separate these muons produced at the primary vertex from those generated in hadron decay or in the interactions with the materials of the detector, \footnote{\url{https://gitlab.cern.ch/atlas/athena/-/blob/main/PhysicsAnalysis/AnalysisCommon/TruthClassification/README.md\#prompt-muons}} prompt muons are those that satisfy \(\text{truthmuon\_IFFType} == 4\), which correspond to the category PromptMuon: it can be only if his \texttt{truthType} is that of an isolated muon and its \texttt{truthOrigin} is a prompt source. Without this constraint the selected sample could be contain also muons coming from heavy hadron decay, produced at distances of the order of centimeters from the primary vertex and classified separately as \(\text{IFFType} == 8\) for B-hadron and \(\text{IFFType} == 9\)  for C-hadron; in-fly-decay muons coming from pions and kaons within the detector, that give rise to a \emph{kink} in the track and they fall into the category of light flavour, \(\text{IFFType} == 10\); muons which truth type, instead, correspond to an electron, \(\text{IFFType} == 12\); and in the end muons which they're tool can't led to any category, or due to missing truth information \(\text{IFFType} == 1\), or because can't be associated to anyone class we juste define, \(\text{IFFType} == 0\), called \texttt{Unkonwn}, which also belong to this category fake muons produced from bad reconstructed traks.

\subsection{Curvature resolution applied on Prompt Muon}

We begin with the simulated $Z\rightarrow\mu\mu$ sample alone, in the ideal configuration. The quantity on which the first part of this analysis is based is the curvature resolution $\mathcal{R}$ defined as the relative difference between the reconstructed curvature and the truth one:

\begin{equation}
  \mathcal{R} \;\equiv\;
  \frac{1/p_{\mathrm{T}}^{\mathrm{reco}} - 1/p_{\mathrm{T}}^{\mathrm{truth}}}
       {1/p_{\mathrm{T}}^{\mathrm{truth}}}
  \;=\;
  \frac{p_{\mathrm{T}}^{\mathrm{truth}}}{p_{\mathrm{T}}^{\mathrm{reco}}} - 1 .
  \label{eq:osservabile}
\end{equation}

\noindent and will be positive when the reconstructed momentum underestimates the truth. We work with curvature resolution and not directly with the momentum because \(1/p_{T}\) is proportional to the sagitta, that is the quantity measured by the detectors \cite{Gluckstern1963}. For each reconstructed muon, we can obtain information about the truth muon by the associated index in the ntupla. One requirement is asked about this muons, that this they have to be \emph{prompt} according to the truth classification \texttt{(truthmuon\_IFFType == 4}). No requirements will be applied to the reconstructed muon. The analysis will be done by comparing both prompt and non-prompt muons, in order to show the effects. The \(R\) distribution is filled in 100 intervals between \(-0.2\) and \(+0.2\), in a inclusive form and in 7 intervals of \(p_T^{truth}\): $0\text{-}20$, $20\text{-}24$, $24\text{-}27$, $20\text{-}24$, $27\text{-}30$, $30\text{-}33$, $33\text{-}37$, $37\text{-}40$, $40\text{-}42$, $42\text{-}44$, $44\text{-}46$, $46\text{-}50$, $50\text{-}60$, $60\text{-}90$ and $90\text{-}200$ GeV. Intervals are widen as the momentum increases to compensate for the low statistics of the sample \(Z\rightarrow\mu^+\mu^-\), passing from tens of millions of muons in the first intervals to hundreds thousand to the latter intervals. Prompt muons make up \(99.25\%\) of the sample\(Z\) in the ideal geometry, leading to a small difference in RMS (\(\simeq1\%\), almost invisible to the naked eye, \cref{fig:res-inclusiva}). 
\begin{figure}[htpb]
  \centering
  \includegraphics[width=\textwidth]{TESI/h_res_prompt_vs_incl_minimal.pdf}
    \caption{Distribution of the curvature resolution $R$ of \cref{eq:osservabile} for muons of the $Z \to \mu^+\mu^-$ sample in the ideal geometry, comparing the inclusive selection with the one restricted to prompt muons (\texttt{truthmuon\_IFFType == 4}). Prompt muons make up $99.25\%$ of the sample, so the two distributions are visually almost indistinguishable and their RMS differ by about $1\%$, $(28.778\pm0.001)\cdot10^{-3}$ against $(28.445\pm0.001)\cdot10^{-3}$. The effect of the non-prompt contamination is not uniform in $p_{\mathrm{T}}$: as shown in \cref{tab:nonprompt_fraction}, it is concentrated in the lowest bin, where non-prompt muons reach $5.3\%$ of the sample and inflate the RMS, while above 20~GeV it drops below $0.1\%$ and the two values coincide to the third decimal digit.}
  \label{fig:res-inclusiva}
\end{figure}
\noindent However, looking at how that \(0.75\%\) of non-prompt muons is distributed per \(p_T^{truth}\) bin, the story is different as shown in \cref{tab:nonprompt_fraction}. Non-prompt contamination is almost focused on the fisrt bin (0-20 GeV, \(5.3\%\)) then start become under \(0.1\%\) and both RMS are egual to the third decimal digit. This is consistent with the fact that prompt muon coming from \(Z\rightarrow\mu^+\mu-\) has a peak in the spectrum around \(m_Z/2\approx45\ GeV\), while non-prompt muon comes from \(\pi/K\) in-flight decay (where the decay happens inside the detector and the trajectory ends up being deflected) or heavy hadron decay has softer spectrum below 20-30 GeV \cite{ATLAS-MUON-EFF-2021}. In the 0-20 GeV isolated bin: the prompt RMS, $(27.289\pm 0.003) \times 10^{-3}$, is practically identical up to 37 GeV bins, showing that the apparent worsening at low pt is not due to the detector resolution but only due to contamination from non-prompts. The quantitative trend is summarized by the RMS of each $p_{\mathrm{T}}^\mathrm{truth}$ interval in \cref{fig:rms-vs-pt}, where we can already see three regimes of the equation \eqref{eq:parametrizzazione}, highly discussed in the next chapter. 
\begin{table}[htbp]
  \centering
  \begin{tabular*}{\textwidth}{@{\extracolsep{\fill}}cccc}
    \toprule
    $p_{\mathrm{T}}^{\mathrm{truth}}$ [GeV] & \% non-prompt & RMS inclusive & RMS prompt \\
    \midrule
    0--20    & 5.30\%   & $29.752 \pm 0.003$ & $27.289 \pm 0.003$ \\
    20--24   & 0.091\%  & $26.770 \pm 0.004$ & $26.761 \pm 0.004$ \\
    24--27   & 0.046\%  & $27.064 \pm 0.004$ & $27.060 \pm 0.004$ \\
    27--30   & 0.028\%  & $27.345 \pm 0.004$ & $27.344 \pm 0.004$ \\
    30--33   & 0.017\%  & $27.646 \pm 0.004$ & $27.645 \pm 0.004$ \\
    33--37   & 0.010\%  & $27.975 \pm 0.003$ & $27.974 \pm 0.003$ \\
    37--40   & 0.007\%  & $28.348 \pm 0.004$ & $28.347 \pm 0.004$ \\
    40--42   & 0.005\%  & $28.579 \pm 0.004$ & $28.578 \pm 0.004$ \\
    42--44   & 0.004\%  & $28.825 \pm 0.004$ & $28.825 \pm 0.004$ \\
    44--46   & 0.003\%  & $29.021 \pm 0.005$ & $29.021 \pm 0.005$ \\
    46--50   & 0.004\%  & $29.287 \pm 0.004$ & $29.287 \pm 0.004$ \\
    50--60   & 0.007\%  & $29.945 \pm 0.004$ & $29.944 \pm 0.004$ \\
    60--90   & 0.008\%  & $31.913 \pm 0.006$ & $31.912 \pm 0.006$ \\
    90--200  & 0.006\%  & $37.260 \pm 0.013$ & $37.260 \pm 0.013$ \\
    \midrule
    0--200   & 0.75\%   & $28.778 \pm 0.001$ & $28.445 \pm 0.001$ \\
    \bottomrule
  \end{tabular*}
  \caption{Fraction of non-prompt muons and RMS of the curvature resolution $R$, inclusive and for prompt muons only, in each $p_{\mathrm{T}}^{\mathrm{truth}}$ bin of the $Z \to \mu^+\mu^-$ sample in the ideal geometry. The uncertainties are obtained with the ROOT method \texttt{GetRMSError}, corresponding to $\sigma_{\mathrm{RMS}} = \mathrm{RMS}/\sqrt{2N}$. The non-prompt contamination is concentrated in the lowest bin, where it reaches $5.3\%$ and visibly inflates the inclusive RMS, while above 20~GeV it falls below $0.1\%$ and the two columns coincide within the quoted uncertainties. The prompt RMS in the first bin, $(27.289\pm 0.003) \times 10^{-3}$, differs by less than $1\%$ from the one of the 20-37~GeV bin, showing that the apparent degradation at low $p_{\mathrm{T}}$ is due to the non-prompt contamination and not to the detector resolution.}
  \label{tab:nonprompt_fraction}
\end{table}
\begin{figure}[htpb]
  \centering
  \includegraphics[width=\textwidth]{TESI/rms_vs_pt_minimal.pdf}
    \caption{RMS of the curvature resolution $R$ as a function of $p_{\mathrm{T}}^{\mathrm{truth}}$ for prompt muons of the $Z \to \mu^+\mu^-$ sample in the ideal geometry, with the values listed in \cref{tab:nonprompt_fraction}. Each point is placed at the mean value of $p_{\mathrm{T}}^{\mathrm{truth}}$ bin, and the horizontal bars indicate the bin width; the vertical bars are the statistical uncertainties $RMS/\sqrt{2N}$ and are too small to be visible. The trend already shows two of the three regimes of \cref{eq:parametrizzazione}, highly discussed in the next chapter.}
  \label{fig:rms-vs-pt}
\end{figure}
\clearpage

\subsection{Pseudorapidity \(\eta\) dependence}

An extension of the analysis resolution in $p_{\mathrm{T}}$ regards how the resolution has a dependence on $|\eta^{\mathrm{truth}}|$. We analyze the same quantity discussed in Section 5.4.1 for prompt muons, but for each $p_{\mathrm{T}}$ we make a plot in 6 $|\eta^{\mathrm{truth}}|$ bins that follows the geometry of the ATLAS detector. For each histogram, we can extrapolate the RMS, resulting in a plateau in $p_{\mathrm{T}}$ as a function of $|\eta^{\mathrm{truth}}|$. The first result shows a peculiar trend for each bin with RMS resolution as a function of $|\eta^{\mathrm{truth}}|$. Horizontal bars are the bin width that follows the detector geometry, while vertical bars are the statistical error on the RSM and are negligible (\cref{fig:RMSeta}). 
\begin{figure}[htbp]
  \centering
  \includegraphics[width=\textwidth]{TESI/plot_RMS_vs_eta.pdf}
    \caption{RMS of the curvature resolution as a function of $|\eta^{\mathrm{truth}}|$ for prompt muons of the $Z \to \mu^+\mu^-$ sample in the ideal geometry, shown separately for the six $p_{\mathrm{T}}^{\mathrm{truth}}$ bins. Horizontal bars indicate the width of each $|\eta|$ bin, chosen to follow the geometry of the detector, while the vertical bars are the statistical uncertainties $\mathrm{RMS}/\sqrt{2N}$ and are too small to be visible; all values are listed in \cref{tab:etavalue}. The resolution is best in the region $0.1 \leq |\eta| < 1.01$ and degrades in the transition region between the barrel and endcap toroids ($1.3 \leq |\eta| < 1.7$), where the bending power of the magnetic field is lowest, and in the innermost bin $|\eta| < 0.1$, where the gap in the spectrometer reserved for the passage of services reduces the coverage. The outermost bin $2.5 \leq |\eta| < 2.8$, beyond the acceptance of the Inner Detector, shows no appreciable degradation with $p_{\mathrm{T}}$.}
  \label{fig:RMSeta}
\end{figure}
We can find a better resolution in the barrel coverage region ($|\eta^{\mathrm{truth}}| \in [0.1-1.01]$), and a worse resolution in the transition area between the barrel and endcap toroids, due to the low and reversal magnetic field ($|\eta^{\mathrm{truth}}| \in [1.3-1.7]$) \cite{ATLAS-MCP-2023}. The last interval, $2.50<|eta|<2.80$, instead does not show any degradation with momentum. Beyond $|\eta|=2.5$, the inner detector coverage is missing, and muons are reconstructed differently \cite{ATLAS-MUON-EFF-2021}. A similar analysis was done for the resolution $R$, then it was reapeated with $\eta^{\mathrm{truth}}$ extrapolating the RMS for each bin of $p_{\mathrm{T}}$, as shown in \cref{fig:etaov}, and for each 36 combination $p_{\mathrm{T}}\times |\eta^{\mathrm{truth}}|$, grouped by $p_{\mathrm{T}}$, all values was reported in \cref{tab:etavalue}, and the RMS errors was evaluated with the relation $RMS/\sqrt{2N}$ where N is the number of enteries.
\begin{figure}[htbp]
  \centering
  \includegraphics[width=\textwidth]{TESI/plot_eta_overlay.pdf}
    \caption{Normalised distributions of the curvature resolution $R = p_{\mathrm{T}}^{\mathrm{truth}}/p_{\mathrm{T}}^{\mathrm{reco}} - 1$ for prompt muons of the $Z \to \mu^+\mu^-$ sample in the ideal geometry, obtained by summing the six $p_{\mathrm{T}}$ bins within each $|\eta^{\mathrm{truth}}|$ interval. Normalization allows the shapes to be compared independently of the number of muons in each interval. The central region $0.10 \leq |\eta| < 1.01$ gives the narrowest and highest peak, while the transition region $1.30 \leq |\eta| < 1.70$ shows the broadest distribution, consistent with the RMS trend of \cref{fig:RMSeta}.}
  \label{fig:etaov}
\end{figure}

\begin{table}[htpb]
    \centering
    \caption{$|\eta|$ RMS for truth muons (IFFType == 4) in different $p_T$ bins}
    \label{tab:etavalue}
    \begin{tabularx}{\textwidth}{l|X|X|X}
    \toprule
    $\eta$ bin & N Entries & RMS & RMS/\sqrt{2N} \\
    \midrule
    \multicolumn{4}{c}{$p_T = 20$--$30$ GeV} \\
    \midrule
    {[}0.00, 0.10) & 1139356 & 0.021453 & \pm\ 0.000014 \\
    {[}0.10, 1.01) & 11378221 & 0.021517 & \pm\ 0.000005 \\
    {[}1.01, 1.30) & 3518384 & 0.025321 & \pm\ 0.000010 \\
    {[}1.30, 1.70) & 4702538 & 0.030539 & \pm\ 0.000010 \\
    {[}1.70, 2.50) & 8481049 & 0.030358 & \pm\ 0.000007 \\
    {[}2.50, 2.80) & 1229333 & 0.040935 & \pm\ 0.000027 \\
    \midrule
    \multicolumn{4}{c}{$p_T = 30$--$40$ GeV} \\
    \midrule
    {[}0.00, 0.10) & 1766564 & 0.023713 & \pm\ 0.000013 \\
    {[}0.10, 1.01) & 16368469 & 0.022795 & \pm\ 0.000004 \\
    {[}1.01, 1.30) & 5064782 & 0.025938 & \pm\ 0.000008 \\
    {[}1.30, 1.70) & 6798602 & 0.031828 & \pm\ 0.000009 \\
    {[}1.70, 2.50) & 12613391 & 0.031164 & \pm\ 0.000006 \\
    {[}2.50, 2.80) & 1773153 & 0.039548 & \pm\ 0.000021 \\
    \midrule
    \multicolumn{4}{c}{$p_T = 40$--$50$ GeV} \\
    \midrule
    {[}0.00, 0.10) & 1681952 & 0.025889 & \pm\ 0.000014 \\
    {[}0.10, 1.01) & 15610727 & 0.023931 & \pm\ 0.000004 \\
    {[}1.01, 1.30) & 4854688 & 0.026618 & \pm\ 0.000009 \\
    {[}1.30, 1.70) & 6542256 & 0.033138 & \pm\ 0.000009 \\
    {[}1.70, 2.50) & 12294293 & 0.031958 & \pm\ 0.000006 \\
    {[}2.50, 2.80) & 1643379 & 0.038563 & \pm\ 0.000022 \\
    \midrule
    \multicolumn{4}{c}{$p_T = 50$--$80$ GeV} \\
    \midrule
    {[}0.00, 0.10) & 716577 & 0.029623 & \pm\ 0.000025 \\
    {[}0.10, 1.01) & 6594322 & 0.025964 & \pm\ 0.000007 \\
    {[}1.01, 1.30) & 2007943 & 0.027924 & \pm\ 0.000014 \\
    {[}1.30, 1.70) & 2647945 & 0.035539 & \pm\ 0.000015 \\
    {[}1.70, 2.50) & 4714557 & 0.033344 & \pm\ 0.000011 \\
    {[}2.50, 2.80) & 604977 & 0.038559 & \pm\ 0.000036 \\
    \midrule
    \multicolumn{4}{c}{$p_T = 80$--$120$ GeV} \\
    \midrule
    {[}0.00, 0.10) & 102400 & 0.038599 & \pm\ 0.000085 \\
    {[}0.10, 1.01) & 941025 & 0.030661 & \pm\ 0.000022 \\
    {[}1.01, 1.30) & 282418 & 0.031617 & \pm\ 0.000042 \\
    {[}1.30, 1.70) & 367259 & 0.041588 & \pm\ 0.000049 \\
    {[}1.70, 2.50) & 622506 & 0.036296 & \pm\ 0.000033 \\
    {[}2.50, 2.80) & 75289 & 0.039425 & \pm\ 0.000103 \\
    \midrule
    \multicolumn{4}{c}{$p_T = 120$--$500$ GeV} \\
    \midrule
    {[}0.00, 0.10) & 30191 & 0.052816 & \pm\ 0.000217 \\
    {[}0.10, 1.01) & 275625 & 0.038423 & \pm\ 0.000052 \\
    {[}1.01, 1.30) & 81725 & 0.038911 & \pm\ 0.000097 \\
    {[}1.30, 1.70) & 104407 & 0.051196 & \pm\ 0.000113 \\
    {[}1.70, 2.50) & 161112 & 0.040641 & \pm\ 0.000072 \\
    {[}2.50, 2.80) & 16947 & 0.041931 & \pm\ 0.000232 \\
    \bottomrule
    \end{tabularx}
\end{table}
\clearpage

\section{Working Point}
For each muon defined in Sections 5.1 and 5.2, quality cuts called Working Point (WP) are applied. This thesis focuses on Prompt Muons, which are associated directly to the primary vertex and are more isolated in the ID and/or calorimeters. Non-prompt Muons, instead, are classified into Light Flavour Muons coming from light hadrons in the secondary vertex, like Pions and Kaons produced by a gluon jet, and Heavy Flavour Muons, that comes from B and C-hadrons. WPs aim to reject light muons, which generally produce bad-quality muons since, having been produced by a secondary decay, still in flight in the detector before becoming muons in the final state, they can cause a fake track that the detector confuses with a prompt muon track. Instead, muons coming from heavy flavour hadrons, like B-hadrons or C-hadrons, produce high-quality muons in the final state that can be easily distinguished from Prompt Muons. WPs are characterised by their purity and efficiency. The efficiency of a prompt muon selection is the probability that one muon that flies through the detector is reconstructed as a real muon that satisfies all the WP criteria. Similarly, the purity of a WP selection is given by one minus the misidentification hadron rate, where this one represents the light hadron fraction reconstructed as a real muon that satisfies the WP criteria. Three WPs are adopted in order to increase purity and decrease efficiency: Loose, Medium and Tight. The Loose WP was designed for the reconstruction of Higgs boson decays in the four-muon final state, which, due to the high muon multiplicity and large signal-to-background ratio, benefits from a higher efficiency at the cost of less purity and larger systematic uncertainties. Medium WP has a good purity-efficiency ratio, maintaining lower systematic uncertainties in both prompt muon identification efficiency and light flavour background cuts, and is suitable for a wide range of analyses. Tight WP selection has the highest purity in terms of light flavour background rejection and applies all the quality cuts used in the Medium WP, along with additional criteria. We will see in Chapter 6 that the resolution will be independent of the choice of WPs.

\subsection{WP criteria}

Identification criteria, noted as \emph{Working Point}, selected muons with respect to the number of measurement stations, fit quality, and compatibility with independent ID and MS measurements \cite{ATLAS-MUON-EFF-2021}. Now it's right to ask whether the resolution depends on the muon WPs quality. In order to do this, were compared the three WPs discussed in Section 5.3: \emph{Loose}, \emph{Medium} and \emph{Tight} (that actually are cumulative, i.e. Tight implies Medium that implies Loose). The three overlayed distribution as shown in \cref{fig:wp-overlay}. The three plateaus are indistinguishable, meaning that the choice of a WPs does not influence the $p_{\mathrm{T}}$ resolution, but only the efficiency.
\begin{figure}[htbp]
  \centering
  \includegraphics[width=\textwidth]{TESI/h_res_wp_overlay.pdf}
    \caption{Normalised distributions of the curvature resolution $R$ of \cref{eq:osservabile} for prompt muons of the $Z \to \mu^+\mu^-$ sample in the ideal geometry, compared for the three identification criteria: Loose, Medium and Tight. The three working points are cumulative, Tight implying Medium and Medium implying Loose. The Loose and Medium curves are indistinguishable and give the same RMS, $0.0282$, while Tight is slightly narrower, $0.0259$, as a consequence of the additional quality requirements it applies. The choice of working point therefore affects the selection efficiency, quantified in \cref{fig:eff_wp}, \cref{fig:eff_eta} and \cref{fig:eff_phi}, rather than the momentum resolution itself.}
  \label{fig:wp-overlay}
\end{figure}
\noindent Since the efficiency is here a WP selection efficiency, defined as the number of muons that satified the WP criteria with respect to the reconstructed muons that are already matched with a truth muon, we obtain $98.7\%$ for Loose, $97.7\%$ for Medium, and $90.4\%$ for Tight, with a constant trend in $p_\mathrm{T}$ (\cref{fig:eff-wp}) and errors evaluated with the Clopper-Pearson method (\emph{TEfficiency}). A uniform binning divided into 50 at 10 GeV each from 0 up to 500 GeV, where Tight increases slightly with $p_{\mathrm{T}}$. Instead, for the pseudorapidity $\eta^{\mathrm{truth}}$, divided into 27 bins in $[-2.7,2.7]$, only for $\eta\simeq0$ Medium decreases to $71\%$ and Tight to $60\%$ \cref{fig:eff-eta}: this is the acceptance hole due to service in the inner barrel region, expected in ATLAS \cite{ATLAS-MUON-EFF-2021}. Loose, instead, is less affected because include segment-tagged and calo-tagged muons, which do not require full coverage of the MS \cite{ATLAS-MUON-EFF-2021}. Same story with respect to the azimuthal coordinate $\phi^{\mathrm{truth}}$, with 32 bins in $[-\pi,\pi]$, which are flat as expected, and no sector hole evident for each WPs (\cref{fig:eff-phi}).
\begin{figure}[htpb]
  \centering
  \includegraphics[width=\textwidth]{TESI/efficiency_vs_pt.pdf}
    \caption{Relative efficiency of the three identification criteria (Loose, Medium and Tight) as a function of the transverse momentum $p_{\mathrm{T}}^{\mathrm{truth}}$, for prompt muons of the $Z \to \mu^+\mu^-$ sample in the ideal geometry, in a uniform binning of $10$~GeV up to $500$~GeV. The efficiency is defined as the fraction of reconstructed muons satisfying each working point with respect to the reconstructed muons matched to a truth muon, and the uncertainties are evaluated with the Clopper--Pearson method. The three criteria are cumulative, Tight implying Medium and Medium implying Loose, and the integrated values are $98.7\%$, $97.7\%$ and $90.4\%$ respectively. All three are essentially flat in $p_{\mathrm{T}}$, with Tight increasing slightly towards higher momenta; the drop visible in the first bin is due to the lower acceptance in the region $|\eta| < 0.1$, caused by the gap in the spectrometer reserved for the passage of services~\cite{ATLAS-MUON-EFF-2021}.}
  \label{fig:eff-wp}
\end{figure}
\begin{figure}[htpb]
  \centering
  \includegraphics[width=\textwidth]{TESI/efficiency_vs_eta.pdf}
    \caption{Relative efficiency of the three identification criteria as a function of the pseudorapidity $\eta^{\mathrm{truth}}$, for prompt muons of the $Z \to \mu^+\mu^-$ sample in the ideal geometry, in 27 bins over $[-2.7, 2.7]$. The efficiency is flat across most of the acceptance, but drops sharply at $\eta \simeq 0$, where Medium falls to $71\%$ and Tight to $60\%$: this is the acceptance hole due to the services in the inner barrel region, expected in ATLAS~\cite{ATLAS-MUON-EFF-2021}. Loose is far less affected, since it also includes segment-tagged and calorimeter-tagged muons, which do not require full coverage of the Muon Spectrometer.}
  \label{fig:eff-eta}
\end{figure}
\begin{figure}[htpb]
  \centering
  \includegraphics[width=\textwidth]{TESI/efficiency_vs_phi.pdf}
    \caption{Relative efficiency of the three identification criteria as a function of the azimuthal coordinate $\phi^{\mathrm{truth}}$, for prompt muons of the $Z \to \mu^+\mu^-$ sample in the ideal geometry, in 32 bins over $[-\pi, \pi]$. All three working points are flat in $\phi$, as expected from the eightfold symmetry of the spectrometer, with no evidence of sector holes; the small modulation visible for Tight follows the alternation between large and small sectors.}
  \label{fig:eff-phi}
\end{figure}

\clearpage

% ==================================================================
%  CAPITOLO 6 - Risoluzione in momento dei muoni e analisi dei residui di disallineamento
% ==================================================================
\chapter{Muon momentum resolution}

\lettrine[lines=3,loversize=0.1,lraise=0.05]{T}{} he precision with which muon momentum is reconstructed directly determines the resolution of the invariant mass of each final state that contains muons, and thus the precision measurement sensitivity and the new physics research based on this channel. In this chapter, the momentum resolution is measured on a simulated sample through the curvature relation discussed in Section 6.2 in a transverse momentum range from the region of the \(Z\) peak to the TeV scale. Then, the measurement is interpreted by a parametrization \cite{ATLAS-MUON-2016, ATLAS-MCP-2023}, discussed in the next sections, which decomposes the sum in quadrature of three terms with different dependence on the transverse momentum of the muon. Since the terms that increase linearly with \(p_T\) denote residual misalignments between the elements of the ATLAS detector, comparing samples produced with ideal geometry and samples produced with the spectrometer deliberately misaligned allows us to isolate and quantify the contribution of residual misalignment to the degradation of high-\(p_T\) resolution. The analysis is developed on \texttt{lxplus} \cite{lxplus}, a CERN interactive login service. Two storage areas were used: personal area AFS \cite{afs} with a storage of \(10 \mathrm{GB}\), used to contain all the Python files and the proxy, and EOS \cite{eos} area with \(1\mathrm{TB}\) storage that, instead, contains all the outputs. The run happens through HTCondor \cite{htcondor} system, a calculus service adopted by CERN. The reason was the reading speed: Python process at a rate of \(10^4\ \mathrm{Hz}\) while HtCondor at \(3\times 10^5 \mathrm{Hz}\), but with more sessions at the same time.



% =====================================================================

% =====================================================================
\section{Parametrization of the resolution}
\label{sec:parametrizzazione}

The trends obtained in the previous section, for each region of $\eta$ and $p_{T}$, are the result of different mechanisms, each dominant in a different kinematic regime. The parameterization of the relative resolution is written as sum in quadrature of three terms\cite{ATLAS-MUON-2016, ATLAS-MCP-2023} :
\begin{equation}
  \frac{\sigma(p_{\mathrm{T}})}{p_{\mathrm{T}}}
  = \frac{r_0}{p_{\mathrm{T}}} \oplus r_1 \oplus r_2 \cdot p_{\mathrm{T}} .
  \label{eq:parametrizzazione}
\end{equation}
The first term accounts for fluctuations in energy loss in the material being crossed and only matters at low momentum; the second, constant, accounts for multiple scattering and inhomogeneities in the description of the local magnetic field; the third, which
grows linearly with $p_{\mathrm{T}}$, describes spatial resolution
position measurements and the residual misalignments between the
detector elements \cite{ATLAS-MCP-2023}. When a charged particle with moentum $p_{\mathrm{T}}$ passes through a magnetic field with intensity $B$ of length $L$ will travel a trajectory whose
Sagitta is worth, within the limit of small deflection
\begin{equation}
  s \simeq \frac{0.3\,B\,L^{2}}{8\,p_{\mathrm{T}}},
  \label{eq:sagitta}
\end{equation}
%
with $s$ and $L$ expressed in meter, $B$ in Tesla and $p_{\mathrm{T}}$ in
GeV \cite{Gluckstern1963}. From the relation
$p_{\mathrm{T}}\propto1/s$, the relative error on Sagitta propagates to the transverse moment with the same percentage value. A $\sigma_s$ uncertainty about Sagitta independent of the moment, as
that produced by the spatial resolution of the detectors and the uncertainty on the position of the chambers, then produces
\begin{equation}
  \frac{\sigma(p_{\mathrm{T}})}{p_{\mathrm{T}}}
  = \frac{\sigma_s}{s}
  = \frac{8\,\sigma_s}{0.3\,B\,L^{2}}\; p_{\mathrm{T}} ,
  \label{eq:sagitta-risoluzione}
\end{equation}
%
that is the third term in the equation 6.2 with $r_2 \simeq 8\sigma_s/(0.3BL^{2})$. The $r_2$ coefficient is then the uncertainty of the Sagitta rewritten in units of $1/p_{\mathrm{T}}$, explaining why the MS misalignment appears in the detectors. MS misalignment produces two effect \cite{ATLAS-ID-ALIGN-2020}: a degradation of the track parameter resolution, which is even more noticeable the higher the momentum, and a parameter $\delta_{\mathrm{sagitta}}$ which alters the curvature in the opposite way for particles of different charge. The first one is the one that appears in $r_2$ and is the focus of this work.


\section{Alignment performance studies}

The $p_{\mathrm{T}}$ resolution analysis described by Eq. 6.1 is now repeated to compare two different configurations of the Muon Spectrometer: one in which all chambers are aligned (\emph{Ideal geometry}) and one in which the chambers are misaligned (\emph{MS misaligned geometry}, discussed in Section 6.5). For this purpose, two simulated samples, $Z$ and $Z'$, produced in both configurations, are used. Only prompt muons with momentum up to $\sim2$~TeV are considered, with the additional selection $\texttt{muon\_Type}==0$, which corresponds to CB muons only. As described in Chapter 5, Combined muons are indeed the only ones with a track in both the ID, which dominates the measurement at low $p_{\mathrm{T}}$, and the MS, which dominates it at high $p_{\mathrm{T}}$, and therefore allow the effect of the chamber misalignment to be observed over the entire momentum range. Dividing the sample into 16 bins of $p_{\mathrm{T}}$ for each of the 6 bins of $|\eta|$, we can extract for each histogram the $68\%$ interquantile range defined in Eq.~6.5:
%
\begin{equation}
  \sigma_{68} = \frac{q_{84} - q_{16}}{2}\,,
  \label{eq:interquantile}
\end{equation}
%
This allow us to consider only the central area of the distribution, using the 16th and 84th percentiles, ignoring what happens in the tails. The error, instead, is just the statistical error $\sigma_{68}/\sqrt{2N}$. For each $|\eta|$ bin we can fit equation 
6.2 considering $\sigma_{68}(p_{\mathrm{T}}^{\mathrm{truth}})/p_{\mathrm{T}}$ performed on a log scale in \cref{fig:resq68-1}. The result of each parameters are shown in \cref{tab:risultati_fit_all}.
\begin{figure}[H]
  \centering
  \includegraphics[width=\textwidth]{TESI/Fit_all_cb_minimal.pdf}
    \caption{Relative momentum resolution $\sigma_{68}/p_{\mathrm{T}}$ as a function of $p_{\mathrm{T}}^{\mathrm{truth}}$ for prompt muons of the combined $Z \to \mu^+\mu^-$ and $Z' \to \mu^+\mu^-$ samples in the ideal (perfectly aligned) geometry. Plots (a), (b), (c), (d), (e) and (f) show separately the six $|\eta|$ bins that follow the detector geometry. For each bin, $\sigma_{68}$ is extracted as half the difference between the 84th and 16th percentiles of the curvature distribution of \cref{eq:osservabile}, and the solid lines are the fits of \cref{eq:parametrizzazione}; the resulting parameters are listed in \cref{tab:risultati_fit_all}.}
  \label{fig:resq68-1}
\end{figure}
The parameter $r_0$ is compatible with zero at low momentum, where the $r_0/p_{\mathrm{T}}$ term carries the most weight, because the CB track is dominated by the ID information. Since there is no material in front of the ID, energy-loss fluctuations are negligible and $r_0$ is therefore consistent with zero. In the region $2.5<|\eta|<2.8$ the ID no longer covers the acceptance (only a few residual CB muons remain), so the reconstruction is governed entirely by the MS. In this case, however, the MS resolution is affected by energy-loss fluctuations in the calorimeter located in front of it, which explains why the value in this region deviates significantly from zero ($0.108\pm0.018$~GeV) and leads to a large $\chi^2/\mathrm{ndf}$ in \cref{fig:resq68-1}~(f). The constant term $r_1$ increases with pseudorapidity, going from a minimum of $0.0139 \pm 0.0022$ in the central region of the detector to a maximum of $0.0256\pm0,0020$ in the transition region of the magnetic field. This trend matches the increase in material crossed and the local magnetic field irregularities that this term accounts for. The observed trend of $r_2$ is consistent with the bending power of the detector's magnet system, as expressed in Eq. 6.3, where $r_2\propto1/(BL^2)$. Indeed, the endcap ($1.7<|\eta|<2.5$), where the bending power is largest, yields a larger sagitta and consequently a small value of $r_2$, $(0.54\pm0.13)\cdot10^{-4}~\text{GeV}^{-1}$. Instead, the transition region, where the bending power is lowest, yields a smaller sagitta and a higher value of $r_2$, $(1.15\pm0.27)\cdot10^{-4}~\text{GeV}^{-1}$. The highest value of $r_2$, however, is observed in the central region ($|\eta|<0.1$), where it reaches $(1.33\pm0.31)\cdot10^{-4}~\text{GeV}^{-1}$. In this case the increase is not related to the bending power but to a geometric effect: a gap in the MS reserved for the passage of services~\cite{ATLAS-MUON-EFF-2021} degrades the measurement of the sagitta, resulting in a high value of $r_2$. The same region also shows the minimum value of the multipl scattering term $r_1$, since the muon traverses the least amount of material. The same effect was already observed in Section 6.1.2 as a degradation of the resolution at high momentum and in Section 6.1.1 as a drop in the efficiency of the Medium and Tight criteria.

\begin{table}[htbp]
    \centering
    \begin{tabular}{c c c c c c}
        \toprule
        Plot & $|\eta^{\mathrm{truth}}|$ & $r_0$ [GeV] & $r_1$ & $r_2$ [$10^{-4}$ GeV$^{-1}$] & $\chi^2/\mathrm{ndf}$ \\
        \midrule
        (a) & $[0.00,\,0.10)$ & $-0.004 \pm 0.039$ & $0.0139 \pm 0.0022$ & $1.33 \pm 0.31$ & $1.79/13$ \\
        (b) & $[0.10,\,1.01)$ & $-0.040 \pm 0.031$ & $0.0174 \pm 0.0017$ & $0.90 \pm 0.22$ & $2.09/13$ \\
        (c) & $[1.01,\,1.30)$ & $\ \ 0.047 \pm 0.008$ & $0.0173 \pm 0.0005$ & $0.89 \pm 0.09$ & $1.08/13$ \\
        (d) & $[1.30,\,1.70)$ & $-0.019 \pm 0.034$ & $0.0256 \pm 0.0020$ & $1.15 \pm 0.27$ & $2.34/13$ \\
        (e) & $[1.70,\,2.50)$ & $-0.024 \pm 0.018$ & $0.0252 \pm 0.0010$ & $0.54 \pm 0.13$ & $1.98/13$ \\
        (f) & $[2.50,\,2.80)$ & $\ \ 0.108 \pm 0.018$ & $0.0205 \pm 0.0009$ & $0.62 \pm 0.12$ & $20.6/13$ \\
        \bottomrule
    \end{tabular}
    \caption{Parameters of the resolution fit of \cref{eq:parametrizzazione} in each $|\eta|$ bin, for the combined $Z$ and $Z'$ samples in the ideal (perfectly aligned) geometry, corresponding to the curves of \cref{fig:resq68-1}. The $\chi^2$ is evaluated with respect to the 13 degrees of freedom: 16 points in $p_{\mathrm{T}}$ minus the 3 free parameters.}
    \label{tab:risultati_fit_all}
\end{table}
\section{Residual Misalignment performance studies}

With the same identical Section 6.4 procedures, we can now analyze the MS misaligned geometry samples of $Z$ and $Z'$, as shown in \cref{fig:q68MS} and the result of each parameters are shown in \cref{tab:risultati_fit_disall}. The values of $r_0$ and $r_1$ do not depend on the alignment of the MS chambers and show no significant variations between the two configurations. What is instead sensitive to a different configuration of the MS chambers is the $r_2$ term, since it depends on the precision of the sagitta measurement, and it is systematically larger than in the ideal configuration, reflecting the introduced misalignment. 
\begin{figure}[htpb]
  \centering
  \includegraphics[width=\textwidth]{TESI/Fit_disall_cb_minimal.pdf}
    \caption{Relative momentum resolution $\sigma_{68}/p_{\mathrm{T}}$ as a function of $p_{\mathrm{T}}^{\mathrm{truth}}$ for prompt muons of the combined $Z \to \mu^+\mu^-$ and $Z' \to \mu^+\mu^-$ samples in the MSmisalign geometry, where the chambers of the middle layer of the Muon Spectrometer are displaced by amounts compatible with the residual precision of the alignment procedure, shown separately for the six $|\eta|$ bins. As in \cref{fig:resq68-1}, $\sigma_{68}$ is extracted from the percentiles of the curvature distribution and the solid lines are the fits of \cref{eq:parametrizzazione}; the resulting parameters are listed in \cref{tab:outputsMSmis}. The overall resolution is degraded with respect to the ideal geometry and the degradation grows with $p_{\mathrm{T}}$, as expected for a term proportional to $p_{\mathrm{T}}$.}
  \label{fig:q68MS}
\end{figure}
\begin{table}[htbp]
    \centering
    \begin{tabular}{c c c c c c}
        \toprule
        Plot & $|\eta^{\mathrm{truth}}|$ & $r_0$ [GeV] & $r_1$ & $r_2$ [$10^{-4}$ GeV$^{-1}$] & $\chi^2/\mathrm{ndf}$ \\
        \midrule
        (a) & $[0.00,\,0.10)$ & $0.017 \pm 0.040$ & $0.0127 \pm 0.0024$ & $2.20 \pm 0.35$ & $0.53/13$ \\
        (b) & $[0.10,\,1.01)$ & $0.014 \pm 0.026$ & $0.0153 \pm 0.0016$ & $1.54 \pm 0.23$ & $0.89/13$ \\
        (c) & $[1.01,\,1.30)$ & $0.069 \pm 0.023$ & $0.0160 \pm 0.0015$ & $1.90 \pm 0.23$ & $1.16/13$ \\
        (d) & $[1.30,\,1.70)$ & $0.056 \pm 0.023$ & $0.0214 \pm 0.0015$ & $2.05 \pm 0.23$ & $2.75/13$ \\
        (e) & $[1.70,\,2.50)$ & $0.070 \pm 0.016$ & $0.0203 \pm 0.0011$ & $1.82 \pm 0.18$ & $2.94/13$ \\
        (f) & $[2.50,\,2.80)$ & $0.217 \pm 0.046$ & $0.0148 \pm 0.0026$ & $2.54 \pm 0.37$ & $9.98/12$ \\
        \bottomrule
    \end{tabular}
    \caption{Parameters of the resolution fit of \cref{eq:parametrizzazione} in each $|\eta|$ bin, for the combined $Z$ and $Z'$ samples in the MS misaligned geometry, corresponding to the curves of \cref{fig:q68MS}. The $\chi^2$ is evaluated with respect to the 13 degrees of freedom: 16 points in $p_{\mathrm{T}}$ minus the 3 free parameters.}
    \label{tab:risultati_fit_disall}
\end{table}
\noindent The residual misalignment term can therefore be extracted by subtracting in quadrature the values of $r_2$ obtained in the configuration with all chambers aligned from those obtained in the configuration with misaligned MS chambers, as shown in Eq. 6.6:
%
\begin{equation}
  r_2^{\mathrm{residual}} = \sqrt{\left(r_2^{\mathrm{MSmisalign}}\right)^2 -
  \left(r_2^{\mathrm{Aligned}}\right)^2}
  \label{eq:residui-disallineamento}
\end{equation}
%
%
The uncertainties on $r_2^{\mathrm{residual}}$ are evaluated by using the central $68\%$ interval of the distribution. This allows us to understand how much alignment impacts the resolution. The results are shown in \cref{tab:result_residual}


\begin{table}[htbp]
  \centering
  \begin{tabular}{clccc}
    \toprule
    & $|\eta^{\mathrm{truth}}|$
    & $r_2^{\mathrm{Aligned}}$ $[10^{-4}\,\mathrm{GeV}^{-1}]$
    & $r_2^{\mathrm{MSmisalign}}$ $[10^{-4}\,\mathrm{GeV}^{-1}]$
    & $r_2^{\mathrm{residual}}$ $[10^{-4}\,\mathrm{GeV}^{-1}]$ \\
    \midrule
    (a) & $[0.00,\,0.10)$ & $1.33 \pm 0.31$ & $2.20 \pm 0.35$ & $1.74^{+0.45}_{-0.59}$ \\
    (b) & $[0.10,\,1.01)$ & $0.90 \pm 0.22$ & $1.54 \pm 0.23$ & $1.26^{+0.29}_{-0.38}$ \\
    (c) & $[1.01,\,1.30)$ & $0.89 \pm 0.09$ & $1.90 \pm 0.23$ & $1.67^{+0.26}_{-0.27}$ \\
    (d) & $[1.30,\,1.70)$ & $1.15 \pm 0.27$ & $2.05 \pm 0.23$ & $1.70^{+0.30}_{-0.38}$ \\
    (e) & $[1.70,\,2.50)$ & $0.54 \pm 0.13$ & $1.82 \pm 0.18$ & $1.74^{+0.19}_{-0.20}$ \\
    (f) & $[2.50,\,2.80)$ & $0.62 \pm 0.12$ & $2.54 \pm 0.37$ & $2.47^{+0.38}_{-0.39}$ \\
    \bottomrule
  \end{tabular}
    \caption{Values of $r_2$ obtained from the fit with all chambers aligned
    ($r_2^{\mathrm{Aligned}}$) and with misaligned MS chambers
    ($r_2^{\mathrm{MSmisalign}}$), and residual misalignment term
    $r_2^{\mathrm{residual}}$ extracted with Eq.~6.6,
    for each $|\eta^{\mathrm{truth}}|$ region. Only combined muons are used.
    The uncertainties on $r_2^{\mathrm{residual}}$ correspond to the $68\%$
    interval of the distribution.}
    \label{tab:result_residual}
\end{table}

\noindent To summarize, the two configurations use the same samples and the same selection, differing only in the alignment of the MS chambers. Any systematic difference can therefore be attributed solely to this misalignment. The $r_2$ term increases in all 6 bins in the MS misaligned configuration with respect to the ideal one. From \cref{tab:result_residual} it can be observed that the residual term is practically constant in bins (a) to (e), compatible within the uncertainties. This is exactly what is expected if the misalignment introduces a uniform contribution to the sagitta measurement, independent of the region. In bin (f) the contribution is higher because, in the absence of the ID, the sagitta is measured only by the MS and only a few CB muons can be adequately reconstructed.


\clearpage


% ==================================================================
%  CAPITOLO 6 - CONCLUSIONI
% ==================================================================
\chapter{Conclusion}
% Testo originale conservato come traccia (era gia' commentato nel file di partenza):
%Muon reconstruction and identification in the ATLAS detector are critical processes to ensure the accuracy of physics analyses, particularly in rejecting background from non-prompt sources. [...]

\begin{todo}[Conclusioni da scrivere]
Da fare: riassumere i risultati ottenuti -- validazione del matching truth/reco,
frazione di muoni prompt (99.25\%), andamento della risoluzione in \(p_T\) e in
\(|\eta|\), parametri \(r_0\), \(r_1\), \(r_2\) e confronto con i campioni disallineati --
e indicare i possibili sviluppi futuri.
\end{todo}


\bibliographystyle{unsrtnat}
\bibliography{references}     % senza estensione .bib
\end{document}